% concatenated from main.tex


\documentclass{amsart}



\usepackage{amsfonts, amssymb}

\usepackage{amssymb,amsmath,amscd,euscript,verbatim,array}
\usepackage{geometry}
\usepackage{anysize}
\usepackage{fancyhdr}
\usepackage{indentfirst}
\usepackage{graphicx}
\usepackage{color}
\usepackage{ifpdf}
\usepackage{marginnote}
% \usepackage{CJK} % removed for arXiv



%% Symbol Definitions --------------------------------------------------------------------



%% Theorem Definitions --------------------------------------------------------------------
%\newcounter{item}[section]
%\newcommand{\thenumber}{\stepcounter{item}\arabic{section}.\arabic{item}}
%\newcommand{\Proof}{\textbf{Proof}\hspace{0.3cm}}
%\newcommand{\End}{\ensuremath{\hfill{\Box}}\\}








%%Kaloshin's Definitions------------------------------------------------------------------
\newcommand{\N}{\mathbb{N}}
\newcommand{\Z}{\mathbb{Z}}
\newcommand{\R}{\mathbb{R}}
\newcommand{\Q}{\mathbb{Q}}
\newcommand{\Lip}{\mathrm{Lip}}
\newcommand{\Id}{\mathrm{Id}}
\newcommand{\ka}{\kappa}
\newcommand{\la}{\lambda}
\newcommand{\al}{\alpha}
\newcommand{\til}{\tilde}
\newcommand{\pa}{\partial}
\renewcommand{\th}{\theta}
\newcommand{\gm}{\gamma}
\newcommand{\ep}{\epsilon}

\newcommand{\Gm}{\Gamma}
\newcommand{\Th}{\Theta}
\newcommand{\A}{\mathbb{A}}
\newcommand{\cA}{\mathcal{A}}
\newcommand{\cM}{\mathcal{M}}
\newcommand{\cT}{\mathcal{T}}
\newcommand{\om}{\omega}
\newcommand{\sg}{\sigma}
\newcommand{\eps}{\varepsilon}
\newcommand{\dt}{\delta}
\newcommand{\lb}{\lambda}
\newcommand{\Lb}{\Lambda}
\newcommand{\cP}{\mathcal P}
\newcommand{\T}{\mathbb{T}}
\newcommand{\disp}{\displaystyle}
\newcommand{\D}[2]{\displaystyle \frac{\partial #1}{\partial #2}}
\newcommand{\Ds}[2]{\frac{\partial #1}{\partial #2}}
\newcommand{\DD}[3]{\displaystyle \frac{\partial^2 #1}{\partial #2\partial #3}}
\newcommand{\DDs}[3]{\frac{\partial^2 #1}{\partial #2\partial #3}}
\newcommand{\DDd}[2]{\displaystyle \frac{\partial^2 #1}{\partial {#2}^2}}
\newcommand{\DDds}[2]{\frac{\partial^2 #1}{\partial {#2}^2}}
\newcommand{\J}{\mathcal{J}}
\newcommand{\twist}{{\mathrm{twist}}}
\newcommand{\low}{{\mathrm{low}}}
\newcommand{\high}{{\mathrm{high}}}
\newcommand{\AM}{{\mathrm{AM}}}
\newcommand{\ovsp}{\mathrm{ovsp}}
\newcommand{\Ss}{\mathbb{S}}



%% Theorem Definitions --------------------------------------------------------------------




\newtheorem{Theorem}{Theorem}
\newtheorem{Definition}{Definition}[section]
\newtheorem{Proposition}[Definition]{Proposition}
\newtheorem{Corollary}[Definition]{Corollary}
\newtheorem{Lemma}[Definition]{Lemma}


\theoremstyle{definition}
\newtheorem{Remark}[Definition]{Remark}




%%Kaloshin's Definitions------------------------------------------------------------------


\newcommand{\bthm}{\begin{Theorem}}
	\newcommand{\ethm}{\end{Theorem}}
\newcommand{\bpr}{\begin{Proposition}}
	\newcommand{\epr}{\end{Proposition}}
\newcommand{\blm}{\begin{Lemma}}
	\newcommand{\elm}{\end{Lemma}}
\newcommand{\bex}{\begin{Exercise}}
	\newcommand{\eex}{\end{Exercise}}
\newcommand{\be}{\begin{equation}}
	\newcommand{\ee}{\end{equation}}
\newcommand{\beal}{\begin{aligned}}
	\newcommand{\enal}{\end{aligned}}
\newcommand{\brm}{\begin{Remark}}
	\newcommand{\erm}{\end{Remark}}

%\fancypagestyle{firstpage}{
%  \fancyhf{} % Çå¿ÕĬÈÏҳüҳ½Å
%  \fancyfoot[L]{2020 Mathematics Subject Classification: 37J40, 37E40}
%  \renewcommand{\headrulewidth}{0pt} % È¥µôҳüºáÏß
%  \renewcommand{\footrulewidth}{0pt} % È¥µôÒ³½ÅºáÏߣ¨¿ÉÉèΪ0.4ptÆôÓã©
%}


\begin{document}


\title[Invariant Graphs for Dissipative Twist Maps]{  On the Dynamics of Invariant Graphs for Dissipative Twist Maps }

\author{Qi Li}
\address{School of Mathematics and Statistics, Beijing Institute of Technology, Beijing 100081, China}
\email{qilicindy@bit.edu.cn}

\author{Lin Wang}
\address{School of Mathematics and Statistics, Beijing Institute of Technology, Beijing 100081, China}
\email{lwang@bit.edu.cn}

\subjclass[2010]{Primary 37J40; Secondary 37E40}


\keywords{Dissipative twist maps, Normally hyperbolicity, Invariant graphs}


\begin{abstract}
For two-parameter families of dissipative twist maps, we investigate the dynamics of invariant graphs as well as the thresholds for their existence and breakdown. Our main results are as follows:
\begin{enumerate}
    \item For arbitrarily small $C^r$ perturbations with $r \geq 1$, invariant graphs with prescribed rotation numbers can be realized by adjusting the parameters;

    \item We characterize sharp perturbations that lead to the complete destruction of all invariant graphs;

    \item When the perturbation fails to be $C^1$, Lipschitz invariant graphs with non-differentiable points may still persist, even though the Lipschitz norm meets the conditions required by the normally hyperbolic invariant manifold theorem.
\end{enumerate}

\end{abstract}



\maketitle



\tableofcontents
	%\newpage
	\section{\sc Introduction}
The investigation of dissipative dynamical systems has been profoundly shaped by the analysis of maps and flows, leading to significant breakthroughs in understanding their structural and behavioral properties. A cornerstone of this field was established by \cite{B}, who in 1932 proved the existence of attractors in dissipative twist maps. Subsequent studies have uncovered a rich spectrum of dynamical phenomena, encompassing periodic solutions, quasi-periodic orbits, KAM tori, and Aubry-Mather sets \cite{CCD1,CCD2,Ca,L1,Ma,MS}.

%Building upon these foundations, \cite{L2} achieved further progress by providing a detailed characterization of the Birkhoff attractor, elucidating its subtle geometric features and their dynamical implications. From a more applied perspective, concrete models have demonstrated the formation of strange attractors in dissipative twist systems, as evidenced by \cite{Za}.

Recent developments have expanded the theoretical framework to higher-dimensional settings. \cite{AHV} successfully generalized the concept of Birkhoff attractors to dissipative maps in higher dimensions, utilizing sophisticated symplectic topological methods for their analysis. Complementary to these advances, rigorous examinations of conformal symplectic systems and their invariant submanifolds have yielded additional insights into these dynamical structures \cite{AA,AF}.
	
	

In this note, we focus on two-parameter families of dissipative twist maps  (see (\ref{Fp}) below) and investigate their dynamical properties under $C^r$ ($r \geq 1$) perturbations. Our main contributions are the following:

\begin{itemize}
    \item {\it  Theorem \ref{M1} and \ref{Mt3}:} We show the diversity of the dynamics on the invariant graphs of perturbed two-parameter families. For any prescribed rotation number, invariant graphs persist under sufficiently small $C^r$  perturbations after appropriate parameter adjustments. However, for a fixed system, the rotation number of the invariant graph  varies under small $C^r$ perturbations in a {\it generic} sense.


    \item {\it  Theorem \ref{M2}, \ref{M22} and \ref{M222}:} We characterize the critical perturbation threshold for the existence and breakdown of invariant graphs in a precise sense. When both normal hyperbolicity and perturbation size decay simultaneously, the persistence or destruction of invariant graphs depends on the ratio of their decay rates. This {\it sharpness} is  quantified by the ratio.

        \item {\it  Theorem \ref{M3}:} While the classical normally hyperbolic invariant manifold (NHIM) theorem \cite[Page 52, Remark 2]{HPS} guarantees that $C^1$ invariant graphs persist under sufficiently small Lipschitz perturbations for $C^1$ normally hyperbolic systems, we construct a  counterexample demonstrating the {\it failure} of this conclusion when both normal hyperbolicity and perturbation size diminish simultaneously.

\end{itemize}





To state the main results in a more precise way, we need to introduce some notions and notations. We denote by $\T:=\R/\Z$.
\begin{Definition}[Dissipative twist map of the cylinder]\label{dissmap}
A {\it dissipative twist map}  of the (infinite) cylinder is a  $C^1$ diffeomorphism $f:\T \times \R \rightarrow \T \times \R  $ that admits a lift
$F:\R^{2}\rightarrow \R^2$, $F(x,y)=(X(x,y),Y(x,y))$, satisfying the following conditions:
\begin{itemize}
    \item[(i)]  ({\it Lift condition}) $F$ is isotopic to the identity;
    \item[(ii)] ({\it Twist condition}) The map $\psi:(x,y)\mapsto(x,X(x,y))$ is a diffeomorphism of $\R^2$;
    \item[(iii)]    ({\it Dissipative condition}) There exist $\lambda_1,\lambda_2\in (0,1)$ such that for all $(x,y)\in \R^2$
\[\lambda_1\leq \det(DF(x,y))\leq \lambda_2.\]
\end{itemize}
\end{Definition}




In this note, we investigate the dynamics of invariant graphs.

\begin{Definition}[$C^0$-invariant graph]\label{c0graph}
$\Gamma\subset\T\times\R$ is called a $C^0$-invariant graph of $f$ if
\begin{itemize}
\item[(i)] $\Gamma=\{(x,\Psi(x)):\; x \in \T\}$, where $\Psi:\T\to\R$ is a continuous function;
\item[(ii)] $\Gamma$ is invariant under the action of  $f$.
\end{itemize}
\end{Definition}

\begin{Remark}
{(i)} By the twist condition, if $f$ is of class $C^1$, then $\Psi$ is a Lipschitz function on $\T$ (see \cite[Proposition 2.2]{H1}).\\
{(ii)} Equivalently, if $F: \R^2 \to \R^2$ denotes a lift of $f$ and $\tilde{\Psi}:\R\to\R$  a $1$ lift of $\Psi$ (which is a $1$-periodic function on $\R$), then  the graph $\tilde{\Gamma}:=\{(x,\tilde{\Psi}(x)) |\ x\in\R\}$  is invariant by $F$.
\end{Remark}




Let $\alpha:=(\alpha_1,\alpha_2)\in \R^2$. Fix $\lambda\in (0,1)$. We consider $F_{\alpha}:=F_{\la,\al}$:
\begin{equation}\label{F1}
F_{\alpha}(x,y):=(x+\alpha_1+\lambda y,\alpha_2+\lambda y).
\end{equation}
The parameter $\lambda$ controls the dissipation, and $\alpha_1,\alpha_2$
 are constants that determine the translation in the $x$ and $y$ directions, respectively. For any parameter pair $(\alpha_1,\alpha_2) \in \mathbb{R}^2$, a direct calculation shows that the integrable system \eqref{F1} possesses a unique invariant graph
\begin{equation}\label{ingm1}
\Gamma = \mathbb{R} \times \left\{\frac{\alpha_2}{1-\lambda}\right\}.
\end{equation}
The map $F_{\alpha}$ restricted on $\Gamma$ reduces to a circle diffeomorphism $g_\al$ with rotation number $\al_1+\frac{\lambda}{1-\lambda}\al_2$. Note that the invariant graph $\Gamma$  is normally hyperbolic (immediately absolutely $r$-normally hyperbolic for any $r\in\N$ in terms of \cite[Definition 2]{HPS}). In fact, for each $(x,y)\in\R^2$,
\[DF(x,y)=\left(
            \begin{array}{cc}
              1 & \la \\
              0 & \la \\
            \end{array}
          \right)=P\left(
            \begin{array}{cc}
              1 & 0 \\
              0 & \la \\
            \end{array}
          \right)P^{-1},
\]
where $P=\left(
            \begin{array}{cc}
              1 & \frac{-\la}{1-\la} \\
              0 & 1 \\
            \end{array}
          \right)$. Hence, we have a $TF$-invariant splitting
\[T\mathbb{T}\times\R|_\Gamma=T\Gamma\oplus N^s.\]
   Moreover, there exists a constant $C$ such that for all $i\in\Z$ and $z\in \Gamma$, there hold
   \begin{equation}\label{hype}
   \frac{1}{C}\la^i|v|\leq |DF^i(z)v|\leq C\la^i|v|,\quad \mathrm{for}\ v\in N^s_z,
   \end{equation}
    \[\frac{1}{C}|v|\leq |DF^i(z)v|\leq C|v|,\quad \mathrm{for}\ v\in T_z\Gamma,\]
    where the constant $C$ can be estimated by the norms of $P$ and $P^{-1}$.

 In \cite{SoW}, the authors consider the following model:
\begin{equation}\label{Fp}
F^\phi_{\alpha}(x,y) = \big(x + \alpha_1 + \lambda y + \phi(x), \alpha_2 + \lambda y + \phi(x)\big),
\end{equation}
where $\phi$ is $1$-periodic and satisfies $\int_{0}^1 \phi(x) dx = 0$.

%We want to consider a perturbation of this family of maps by a $C^1$  function $\phi:\R\to\R$ which is $1$-periodic and satisfies $\int_{0}^1\phi(x)dx=0$, namely
%\begin{equation}\label{Fp}
%F^{\phi}_{\al}(x,y)=(x+\alpha_1+\lambda y,\alpha_2+\lambda y+\phi(x+\alpha_1+\lambda y)).
%\end{equation}
Clearly, $\det DF_{\al}^\phi(x,y)=\lambda$ for all $(x,y)\in \R^2$, hence this induces a family of dissipative twist maps. The invariant graphs of integrable two-parameter families of dissipative twist maps exhibit normal hyperbolicity. Their persistence under $C^r$ perturbations follows from the NHIM theorem  \cite{HPS}. For $\lambda \in (0,1)$, if an invariant graph exists, then it serves as the \emph{global attractor} of the system. Analogously, one can extend the corresponding results obtained in this note to the opposite case with $\lambda>1$ and obtain a dynamical characterization of invariant graphs with repelling properties.\\




If $\lambda=1$ and $\alpha=0$, then $F^{\phi}_{\lambda,\alpha}$ reduces to an exact area-preserving twist map. Specifically,
\begin{equation}\label{Fpar}
F^{\phi}(x,y) = (x + y+\phi(x), y + \phi(x)).
\end{equation}

The existence of an invariant graph for $F^{\phi}$ and the dynamics on this graph are closely related to the $C^r$ topology in which the perturbation is considered. When restricted to the invariant graph, $F^{\phi}$ becomes an orientation-preserving circle homeomorphism, and its dynamics are essentially determined by the rotation number's arithmetic properties. Furthermore, based on the {Denjoy--Herman--Yoccoz} theory~\cite{De,H11,Y1}, we can analyze finer dynamical properties on the invariant graphs, such as whether the system is conjugate to a rigid rotation, or even the precise sense in which such conjugacy holds.

 \vspace{1em}

\subsection{The dynamics on invariant graphs}
Recall that a number $\omega \in \mathbb{R}$ is called \emph{Diophantine} if there exist a constant \( D > 0 \), and an exponent \( \tau \geq 0 \), such that for all   $p,q\in \Z$ and $q\neq 0$,
\begin{equation}\label{diop}
 \left|\omega - \frac{p}{q}\right| \geq \frac{D}{|q|^{2+\tau}}.
 \end{equation}
The exponent $\tau$ is called the \emph{Diophantine exponent}. If $\tau = 0$, we say $\omega$ is of \emph{constant type}. Irrational numbers that are not Diophantine are called \emph{Liouville numbers}.

Herman \cite{H33} proved that if $\|\phi\|_{C^3}$ is sufficiently small, then the invariant graph with constant-type rotation number persists, and the restricted dynamics on it are $C^1$-conjugate to a rigid rotation. On the other hand, for any $\tau_0 > 0$, one can construct a perturbation $\phi_0$ that remains arbitrarily small in the $C^3$ topology but destroys the invariant graph with Diophantine exponent $\tau_0$ (see \cite{H1,W}).

Furthermore, Mather \cite{M4} showed that each invariant graph with Liouville rotation number can be destroyed by $C^\infty$-small perturbations. Forni \cite{F} later proved that for a special subclass of frequencies (a proper subset of non-Brjuno numbers), even $C^\omega$ (real-analytic) perturbations suffice to break the invariant graph.

For $\lambda \in (0,1)$, the system becomes dissipative. The unperturbed system admits a unique invariant graph, which is normally hyperbolic. Thus, the persistence of the invariant graph under perturbation is guaranteed by the NHIM theorem, provided the perturbation is sufficiently small in the $C^r$ topology ($r\geq 1$). A natural question arises:

\vspace{1ex}


\begin{itemize}
\item \textbf{Question 1:} {\it If the perturbation is  required to be $C^r$-small, how general can the dynamics on the persisted invariant graph be? In particular, can the rotation number be arbitrary?}
\end{itemize}

\vspace{1ex}


 In the conservative case, the answer is negative---for instance, under $C^3$-small perturbations, only invariant graphs with constant-type rotation numbers persist generally. This question serves as a key motivation for our investigation. We will demonstrate (see Theorem \ref{M1} below) that the answer is positive in the dissipative setting if we consider the family of $\{F^{\phi}_{(\al_1,\al_2)}\}_{(\al_1,\al_2)\in\R^2}$ with two parameters.




  In the case with $\al_2=0$, $F^{\phi}_{\al}$ is also referred as to an exact conformally symplecitc twist map.
\begin{Theorem}\label{M1}
Given $r\in [1,+\infty)$ and $0<\eps\ll 1$, there exists $\delta=O(\eps)$, such that if the perturbation $\phi$ satisfies $\|\phi\|_{C^r}\leq \delta$, then the following statements hold true.
\begin{itemize}
\item [(1)] The perturbed map $F^{\phi}_{\al}$ still admits a $C^r$-invariant graph $\tilde{\Gamma}:=\{(x,\tilde{\Psi}_{\al}(x))\ |\ x\in \R\}$. Moreover,
\[F^{\phi}_{\al}(x,\tilde{\Psi}_{\al}(x))=(g_{\al}(x), \tilde{\Psi}_{\al}(g_{\al}(x))),\]
where $g_{\al}(x):=x+\al_1+\lambda\tilde{\Psi}_{\al}(x)+\phi(x)$.
\item [(2)] For each $x\in \R$, the function $\tilde{\Psi}_\alpha(x)$ is uniformly Lipschitz with respect to the parameter $\alpha:=(\alpha_1,\alpha_2)\in \R^2$ and the Lipschitz constant is also independent of $x$.
\item [(3)] Denote $\rho(g_{\al})$ to be the rotation number of $g_{\al}$.  Given $\al_2\in\R$ (resp. $\al_1\in\R$), for each $\omega\in\R$, there exists $\bar{\al}_1\in\R$ (resp. $\bar{\al}_2\in\R$) such that $\rho(g_{(\bar{\al}_1,\al_2)})=\omega$ (resp. $\rho(g_{({\al}_1,\bar{\al}_2)})=\omega$).
\end{itemize}
\end{Theorem}


\begin{Remark}
According to Theorem~\ref{M1}, the two-parameter family $F_{\al}$ defined by~\eqref{F1} preserves invariant graphs with arbitrary rotation numbers  (up to small deformations) under small perturbations in the $C^r$ ($r \geq 1$) topology. The regularity of these invariant graphs is entirely determined by the perturbation. From the construction of $F^{\phi}_{\alpha}$, the orientation-preserving circle homeomorphism induced by the perturbed map $\left.F^{\phi}_{\alpha}\right|_{\tilde{\Gamma}}$ restricted to the invariant graph $\tilde{\Gamma}$ shares the same regularity as the invariant graph itself. Consequently, classical circle map theory  can be also applied to obtain more refined dynamical properties on the invariant graph.
\end{Remark}

 According to the NHIM theorem, when subjected to sufficiently small $C^1$ perturbations, the system \eqref{F1} retains a unique perturbed invariant graph. While the NHIM theorem guarantees the persistence of an invariant graph under perturbation, it does not preserve the rotation number generically.
%More precisely, we fix  $\alpha_2 = 0$, and write $F^\phi_{\alpha_1} := F^\phi_{(\alpha_1,0)}$ for simplicity.
\begin{Theorem}\label{Mt3}
Given $r \in [1,+\infty)$ and a generic $\alpha_1 \in \mathbb{R}$, there exists a sequence of $C^\infty$ functions $\{\psi_n\}_{n \in \mathbb{N}}$ with $\|\psi_n\|_{C^r} \to 0$ as $n \to +\infty$ such that:
\begin{itemize}
    \item $F^{\psi_n}_{(\alpha_1,0)}$ admits an invariant graph $\tilde{\Gamma}_n$;
    \item The induced circle map on $\tilde{\Gamma}_n$ has a rational rotation number different from $\alpha_1\in\R\backslash\mathbb{Q}$.
\end{itemize}
\end{Theorem}



 \vspace{1em}

\subsection{The threshold for the existence of invariant graphs}


When the perturbation exceeds a certain threshold, this invariant graph undergoes breakdown. This naturally leads to the  question of determining the critical perturbation strength at which invariant graphs are either preserved or destroyed.

Drawing an analogy with conservative systems - where similar questions are addressed by KAM theory and converse KAM theory - we may investigate two distinct problems:
\begin{enumerate}
    \item [(1)] The critical perturbation required to destroy an invariant graph with prescribed rotation number
    \item [(2)] The critical perturbation needed to destroy all possible invariant graphs
\end{enumerate}

In view of Theorem \ref{Mt3},
 the second problem concerning the destruction of all invariant graphs is more natural to consider for the two-parameter family $\{F_{\al}\}_{\alpha\in\mathbb{R}^2}$. Specifically, we formulate the following question:

\vspace{1ex}

\begin{itemize}
\item \textbf{Question 2:} {\it What is the sharp perturbation required to cause the breakdown of all invariant graphs in the two-parameter family $\{F_{\al}\}_{\alpha\in\mathbb{R}^2}$?}
\end{itemize}

\vspace{1ex}

By the NHIM theorem, the perturbation in Question 2 cannot be arbitrarily small in the $C^1$ topology.   The main result obtained in \cite{SoW} can be stated as follows:




\begin{Proposition}\label{Mt}
Given $\lambda \in (0,1)$ and $0 < \varepsilon \ll 1$, there exists a sequence of trigonometric polynomials $\{\varphi^{\lambda}_n\}_{n \in \mathbb{N}}$ such that all $C^0$-invariant graphs for the family $\{F_{\al}\}_{\alpha \in \mathbb{R}^2}$ can be destroyed by perturbing the maps with $\{\varphi^{\lambda}_n\}_{n \in \mathbb{N}}$. Moreover, the following properties hold:
\begin{itemize}
    \item[(a)] $\|\varphi^{\lambda}_n\|_{C^{1 - \varepsilon}} = O\left(\frac{1}{n^\varepsilon}\right)$ as $n \to \infty$;
    \item[(b)] For every $n \in \mathbb{N}$, $\|\varphi^{\lambda}_n\|_{C^1} \leq 1$ for all $\lambda \in (0,1)$, and $\|\varphi^{\lambda}_n\|_{C^1} = O(1 - \lambda)$ as $\lambda \to 1^-$.
\end{itemize}
\end{Proposition}

Item~(a) in Proposition~\ref{Mt} is sharp, as the NHIM theorem ensures the persistence of invariant graphs under small $C^1$ perturbations. We will show that Item~(b) is  also sharp  to cause the complete breakdown of all invariant graphs. To illustrate what we mean by ``sharp," we consider the following question:

\vspace{1ex}

\begin{itemize}
\item \textbf{Question 3:} {\it  For all perturbations satisfying
\[
\|\phi^{\lambda}\|_{C^1}=O((1-\lambda)^{\gamma}), \quad \text{as} \quad \lambda\to 1^-,
\]
if there exists $\tilde{\al}\in\R^2$ such that $F^{\phi^\lambda}_{\tilde{\al}}$ admits an invariant graph, what is the infimum of all possible values of $\gamma$?}
\end{itemize}

\vspace{1ex}





Proposition~\ref{Mt} establishes that $\gamma \geq 1$ is necessary. We will prove the upper bound $\gamma \leq 2$.  Let $\varphi$ be a Lipschitz function on $\T$. Recall the Lipschitz semi-norm defined by:
\[\|\varphi\|_{\mathrm{Lip}}:=\sup_{x\neq y}\frac{|\varphi(x)-\varphi(y)|}{|x-y|},\]
which is also referred to as the Lipschitz constant of $\varphi$. If $\varphi$ is $1$-periodic and satisfies $\int_{0}^1 \phi(x) dx = 0$. By the mean value theorem, we have $\|\varphi\|_{C^0}\leq \|\varphi\|_{\mathrm{Lip}}$. For each $\al_1\in\R$, we choose  $\tilde{\alpha} = (\alpha_1, 0) \in \mathbb{R}^2$. Our main quantitative result for model (\ref{Fp}) is the following:

\begin{Theorem}\label{M2}
Given $\lambda \in (0,1)$, for each $\al_1\in\R$ and  Lipschitz perturbation $\phi^\lambda$ satisfying
\begin{equation}\label{contrac}
\|\phi^\lambda\|_{\mathrm{Lip}} < (1-\sqrt{\lambda})^2,
\end{equation}
the map $F^{\phi^\lambda}_{(\al_1,0)}$ admits a unique Lipschitz invariant graph.
\end{Theorem}

\begin{Remark}
For notational clarity, we adopt the following conventions throughout:
\begin{itemize}
\item $u \lesssim v$ (resp. $u \gtrsim v$) denotes $u \leq Cv$ (resp. $u \geq Cv$)
\item $u \sim v$ means $\frac{1}{C}v \leq u \leq Cv$
\end{itemize}
for some positive constant $C$.
It is easy to see that $(1-\sqrt{\lambda})^2\sim (1-\lambda)^2$ as $\lambda\to 1^-$. For specific perturbations like the dissipative standard map
\begin{equation}\label{contrac33}
\varphi(x) = \frac{\kappa}{2\pi} \sin(2\pi x),
\end{equation}
Bohr \cite{Bo} showed that invariant graphs are destroyed when $\kappa >  \frac{2(1+\lambda)}{2+\lambda}$.  Fixing $\alpha = (\alpha_1, 0) \in \mathbb{R}^2$, we obtain the existence  of invariant graphs when $\kappa \leq (1-\sqrt{\lambda})^2$. It follows that Theorem \ref{M2} is optimal in the sense:
\[\Delta(\lambda):=\frac{2(1+\lambda)}{2+\lambda}-(1-\sqrt{\lambda})^2\to 0^+,\quad \text{as} \quad \lambda\to 0^+.\]


For comparison, in the conservative case, current theoretical results (to the best of our knowledge) show:
\begin{itemize}
\item Complete breakdown occurs when $\kappa > \frac{4}{3}$ \cite{M1};
\item Existence persists for rotation number $\frac{1+\sqrt{5}}{2}$ when $\kappa \leq 0.029$ \cite[Page 197]{H33}.
\end{itemize}
\end{Remark}



If the perturbation $\phi^\lambda$ is a $C^1$ function, then under condition (\ref{contrac}) of Theorem \ref{M2},  the regularity of the invariant graph can be improved to $C^1$. This regularity enhancement is based on both the contractive property (hyperbolicity in more general cases) of the mapping $F^{\phi^\lambda}_{(\al_1,0)}$ and the invariance of the graph with respect to $F^{\phi^\lambda}_{(\al_1,0)}$.





\begin{Theorem}\label{M22}
Given $\lambda \in (0,1)$, for each $C^1$ perturbation $\phi^\lambda$ satisfying
\[
\|\phi^\lambda\|_{C^1} < (1-\sqrt{\lambda})^2,
\]
the map $F^{\phi^\lambda}_{(\al_1,0)}$ admits a unique $C^1$ invariant graph.
\end{Theorem}


In view of (\ref{hype}), the parameter $\lambda$ quantifies the degree of normal hyperbolicity: as $\lambda \to 1^-$, the normal hyperbolicity weakens, while as $\lambda \to 0^+$, it strengthens.
Theorem \ref{M22} presents a quantitative version of the NHIM theorem, establishing the relationship between the existence of invariant graphs and the ratio of decay rates as both normal hyperbolicity and the $C^1$-norm of the perturbation  decrease simultaneously.

 \vspace{1em}

Inspired by \cite{Ma} and employing normal form techniques, we perform a sequence of coordinate changes to obtain a more refined result.


\begin{Theorem}\label{M222}
Let $\gamma \in (1,2)$ and let $\alpha_1$ be a Diophantine number. Then there exists $\lambda_0 \in (0,1)$ such that for all $\lambda \in [\lambda_0, 1)$ and for each real-analytic perturbation $\phi^\lambda \in C^\omega(\mathbb{T})$ satisfying
\[
\|\phi^\lambda\|_{C^\omega} < (1 - \lambda)^{\gamma},
\]
the map $F^{\phi^\lambda}_{(\alpha_1, 0)}$ admits a unique $C^1$ invariant graph, where $\|\cdot\|_{C^\omega}$ is defined by~\eqref{renorm}.
\end{Theorem}

Define for $r\in [1,+\infty]\cup\{\omega\}$,
\[
\mathcal{S}^r_\lambda := \left\{ \phi^\lambda \in C^r(\mathbb{T})\ \middle|\ \exists\, \tilde{\alpha} \in \mathbb{R}^2 \text{ such that } F^{\phi^\lambda}_{\tilde{\alpha}} \text{ admits a } C^1 \text{ invariant graph} \right\},
\]
\[
\Lambda^r:= \left\{ \gamma \in \mathbb{R}\ \middle|\ \phi^\lambda \in \mathcal{S}^r_\lambda,\ \|\phi^\lambda\|_{C^r} = O((1 - \lambda)^{\gamma}) \text{ as } \lambda \to 1^- \right\}.
\]

Combining Proposition~\ref{Mt} and Theorem~\ref{M222}, we conclude that
\[
\inf \Lambda^\omega = 1,
\]
and the infimum is not attained. This provides a complete answer to Question 3 in the real-analytic category. However, based on Proposition~\ref{Mt} and Theorem~\ref{M22}, we only know that in the $C^1$ category,
\[
1 \leq \inf  \Lambda^1 \leq 2.
\]


\subsection{The regularity of invariant graphs}


Based on Theorem \ref{M2} and Theorem \ref{M22}, we are naturally led to the following question:

\vspace{1ex}

\begin{itemize}
\item \textbf{Question 4:} {\it When the perturbation is only Lipschitz continuous (with non-differentiable points) but satisfies (\ref{contrac}), what regularity does the invariant graph possess?}
\end{itemize}

\vspace{1ex}


%In general, the answer to Question 3 is negative. The twist condition alone can only guarantee that the preserved invariant graph remains Lipschitz continuous. We demonstrate this by constructing appropriate perturbations of system (\ref{Fp}).

We construct an example to demonstrate the existence of a Lipschitz perturbation satisfying condition (\ref{contrac}) for which the preserved Lipschitz invariant graph contains non-differentiable points. It shows that Theorems \ref{M2} and \ref{M22} are mutually non-implicative.
\begin{Theorem}\label{M3}
Given $\lambda \in (0,1)$, $0 < \varepsilon \ll 1$, and an irrational number $\omega$, there exists a sequence of Lipschitz functions $\{\phi_n^{\lambda}\}_{n \in \mathbb{N}}$ and a sequence $\{\bar{\alpha}_n\}_{n \in \mathbb{N}} \subset \mathbb{R}$ such that the map $F^{\phi_n^\lambda}_{(\bar{\alpha}_n,0)}$ admits a Lipschitz invariant graph $\tilde{\Gamma}_n$ that contains non-differentiable points. Moreover:
\begin{itemize}
    \item[(I)] $\|\phi^{\lambda}_n\|_{C^{1-\varepsilon}} = O\left(\frac{1}{n^\varepsilon}\right)$ as $n \to \infty$;
    \item[(II)] For all $n \in \mathbb{N}$, $\|\phi^{\lambda}_n\|_{\mathrm{Lip}} \leq (1 - \sqrt{\lambda})^2$;
    \item[(III)] The circle map induced by the restriction of $F^{\phi_n^\lambda}_{(\bar{\alpha}_n,0)}$ to $\tilde{\Gamma}_n$ has rotation number $\omega$.
\end{itemize}

\end{Theorem}


%\begin{Remark}
%The preserved invariant graphs in Theorem \ref{M2}  are graphs of Lipschitz functions which have non-differentiable points, which consequently cannot be detected by the NHIM theorem.
%In \cite{Ha}, Hall constructed a dissipative $C^\infty$ diffeomorphism that preserves such Lipschitz graphs which have non-differentiable points. In his example, the rotation number of the invariant graph is well-approximable by rationals (e.g., Liouville-type numbers). Hall  raised the question of whether such examples exist for arbitrary irrational rotation numbers \cite[Page 117, Remark (1)]{Ha}.
%Theorem \ref{M2} provides an affirmative answer for the $C^1$ case: for any prescribed irrational rotation number $\omega$, we can construct such examples within the class of $C^1$ mappings.
%\end{Remark}

\begin{Remark}
According to \cite[Page 52, Remark 2]{HPS}, for a $C^1$ normally hyperbolic map $F$, a $C^1$ invariant graph of $F^\phi$ persists under a Lipschitz-small perturbation $\phi$. However, Theorem \ref{M3} demonstrates that this persistence fails when both the perturbation size and the normal hyperbolicity decay simultaneously.
\end{Remark}

\begin{Remark}
While dissipative twist maps can, in general, admit invariant curves that are not graphs (see Proposition 15.3 in~\cite{L2}), the perturbation constructed in~\cite{L2} is not small in the $C^0$ topology. Thus, the problem of finding critical perturbations that break all invariant curves of the two-parameter family $\{F_{\al}\}_{\alpha\in\mathbb{R}^2}$ (regardless of whether they are graphs) remains open.
\end{Remark}


\subsection*{Organization of the note}
The note is organized as follows. In Section \ref{S3}, we investigate the diversity of dynamics on invariant graphs. By adjusting either of two parameters, we show that invariant graphs can exhibit arbitrary frequencies, thereby proving Theorem \ref{M1} and \ref{Mt3}.
In Section \ref{S4}, we examine two-parameter families of dissipative twist maps, focusing on the critical perturbation size that leads to the breakdown of all invariant graphs. According to the NHIM theorem, for any given $\lambda \in (0,1)$, the perturbation required to destroy invariant graphs cannot be arbitrarily small in the $C^1$-norm (or more generally, in the Lipschitz semi-norm). We prove three quantitative versions of the NHIM theorem in our setting ( see Theorem \ref{M2}, \ref{M22} and \ref{M222}). As $\lambda \to 1^-$, both the perturbation size and normal hyperbolicity decay simultaneously. The {\it almost sharpness} discussed in this section is determined by the ratio of their decay rates.
In Section \ref{S5}, we construct an example (Theorem \ref{M3})  demonstrating that Theorems \ref{M2} and \ref{M22} are mutually non-implicative. This result complements the classical NHIM theorem by showing that for $C^1$ systems under sufficiently small Lipschitz perturbations, invariant graphs need not remain $C^1$ when both the perturbation size and normal hyperbolicity diminish simultaneously.





	 \vspace{2em}

 \noindent\textbf{Acknowledgement.}
The authors would like to thank Jessica Massetti for her patient and helpful explanations regarding \cite{Ma}. This work was partially supported by the National Natural Science Foundation of China (Grant No.~12122109).

	
	
	
	 \vspace{2em}

		
%	\section{\sc Preliminaries}\label{S2}
%
%	 \vspace{2em}



\section{\sc The dynamics on invariant graphs}\label{S3}

\subsection{Persistence of invariant graphs}
This subsection is devoted to the proof of Item~(1) in Theorem~\ref{M1}. We begin by recalling some foundational concepts and establishing key notations.

Following Herman, we denote by $\mathrm{Diff}^r_+(\mathbb{R})$ (resp. $\mathrm{Diff}^r_+(\mathbb{T})$) the group of orientation-preserving $C^r$ diffeomorphisms on $\mathbb{R}$ (resp. $\mathbb{T}$), where $r \in [0,+\infty) \cup \{+\infty\} \cup \{\omega\}$. The universal covering space of $\mathrm{Diff}^r_+(\mathbb{T})$ can be represented as
\[
D^r(\mathbb{T}) := \{f \in \mathrm{Diff}^r_+(\mathbb{R}) \mid f - \mathrm{Id} \in C^r(\mathbb{T})\}.
\]
For any $f \in D^r(\mathbb{T})$, the rotation number $\rho(f)$ is well-defined. By \cite[Proposition 2.1]{SoW}, we have

\begin{Proposition}\label{hf1}
The map $F^\phi_{\al}$ admits a $C^0$-invariant graph $\tilde{\Gamma} := \{(x, \tilde{\Psi}_{\alpha}(x)) \mid x \in \mathbb{R}\}$ if and only if there exists $g_{\alpha} \in D^r(\mathbb{T})$ satisfying the functional equation
\begin{equation}\tag{A}\label{eq:main}
\frac{1}{1+\lambda}g_{\alpha}(x) + \frac{\lambda}{1+\lambda}g_{\alpha}^{-1}(x) = x + \frac{1}{1+\lambda}\left((1-\lambda)\alpha_1 + \lambda\alpha_2 + \phi(x)\right) \quad \forall x \in \mathbb{R}.
\end{equation}
\end{Proposition}

The invariance of $\tilde{\Gamma}$ implies the  relation
\begin{equation}\label{finvar}
F^\phi_{\al}(x, \tilde{\Psi}_{\alpha}(x)) = (g_{\alpha}(x), \tilde{\Psi}_{\alpha}(g_{\alpha}(x))),
\end{equation}
from which we immediately deduce that
\begin{equation}\label{xw2}
g_{\alpha}(x) = x + \alpha_1 + \lambda\tilde{\Psi}_{\alpha}(x)+\phi(x),
\end{equation}
\begin{equation}\label{psss}
\phi(x)=\tilde{\Psi}_\al(g(x))-\lambda\tilde{\Psi}_\al(x)-\al_2.
\end{equation}
Let $\pi_1: (x,y) \mapsto x$ and $\pi_2: (x,y) \mapsto y$ denote the canonical projections. For fixed $\lambda \in (0,1)$, we simplify notation by writing $F^{\phi}_\alpha(x,y) := F^{\phi}_{\al}(x,y)$. The following persistence result follows from the NHIM (Normally Hyperbolic Invariant Manifold) theorem:

\begin{Proposition}\label{nhim}
For any $r \in [1,+\infty)$ and $0 < \eps \ll 1$, there exists $\delta = O(\eps)$ such that if $\|\phi\|_{C^r} \leq \delta$, then the perturbed map $F^{\phi}_\alpha$ admits a $C^r$-invariant graph $\tilde{\Gamma}_\alpha^\phi := \{(x, \tilde{\Psi}_{\alpha}(x)) \mid x \in \mathbb{R}\}$ with the estimate
\[
\left\|\tilde{\Psi}_\alpha - \frac{\alpha_2}{1-\lambda}\right\|_{C^r} \leq \eps.
\]
\end{Proposition}



\subsection{Lipschitz dependence of invariant graphs on parameters}
We now turn to the verification of Item (2) in Theorem \ref{M1}.  We will use the following notion from \cite{Co}.
\begin{Definition}
Let $X$ be a metric space, $\mathcal{C}$ a compact subset of $X$ and $f:X\to X$ a continuous map, then we call $\mathcal{C}$ an attractor block for $f$ if
\[f(\mathcal{C})\subseteq \mathrm{interior}\ \mathcal{C}.\]
The set
\[\cap_{i=1}^\infty f^i(\mathcal{C})\subseteq \mathcal{C},\]
which is the largest invariant set in $\mathcal{C}$, is called the attractor for the attractor block.
\end{Definition}

Without ambiguity, We  use $\Psi_\al$ to denote $\tilde{\Psi}_\al$ for simplicity.
\begin{Proposition}\label{lipo}
 If $\|\phi\|_{C^1}\leq \delta_0$ where $\delta$ is the constant determined by Proposition \ref{nhim}, then the function $\tilde{\Psi}_\alpha(x)$ is uniformly Lipschitz with respect to the parameter $\alpha:=(\alpha_1,\alpha_2)\in \R^2$. More precisely, for each $x\in \R$ and $\al,{\al'}\in \R$,
\[|\tilde{\Psi}_\al(x)-\tilde{\Psi}_{\al'}(x)|\leq \frac{1}{1-\lambda}|\al_2-\al'_2|.\]
\end{Proposition}
\begin{proof}
Fix the perturbation $\phi:\T\to\R$. For simplicity, we denote $F_\alpha := F^{\phi}_\alpha$. Given $\alpha' := (\alpha'_1, \alpha'_2) \in \R^2$, we consider the transformation $\Phi:\R^2\to\R^2$ defined by
\[
\Phi(x,y) = (x, y - \Psi_{\alpha'}(x)).
\]
Its inverse is given by
\[
\Phi^{-1}(x,y) = (x, y + \Psi_{\alpha'}(x)).
\]
Define the conjugated map
\[
K_{\alpha}(x,y) := \Phi \circ F_\alpha \circ \Phi^{-1}(x,y).
\]
A direct calculation yields
\begin{align*}
K_{\alpha}(x,y) = \biggl(x +  \alpha_1+\lambda(y + \Psi_{\alpha'}(x)) +\phi(x), \lambda(y + \Psi_{\alpha'}(x)) + \alpha_2 + \phi(x) - \Psi_{\alpha'}(x)\biggr).
\end{align*}

Let $\chi(x,y,\alpha) := \pi_2 K_\alpha(x,y)$. Since $\Psi_{\alpha'}$ is an invariant graph for $F_{\alpha'}$, Proposition~\ref{hf1} implies
\[
F_{\alpha'}(x, \Psi_{\alpha'}(x)) = (g_{\alpha'}(x), \Psi_{\alpha'}(g_{\alpha'}(x))),
\]
and consequently $\chi(x,0,\alpha') = 0$ for all $x \in \R$. Define
\[
G(x,y,\alpha) := \chi(x,y,\alpha) - y.
\]
The partial derivatives satisfy:
\begin{align}
\frac{\partial \chi}{\partial y}(x,y,\alpha) = \lambda,\quad \frac{\partial \chi}{\partial \alpha_1}(x,y,\alpha) = 0,\quad \frac{\partial \chi}{\partial \alpha_2}(x,y,\alpha) = 1.
\end{align}

It follows that
\begin{equation}\label{contrpo}
0 < \frac{\partial \chi}{\partial y}(x,y,\alpha)  < 1
\end{equation}
for all $x \in \R$.

Applying the implicit function theorem, there exists a neighborhood $U(\alpha')$ and a $C^1$ function $\sigma:\R \times U(\alpha') \to \R$ such that
\[
\chi(x, \sigma(x, \alpha), \alpha) = \sigma(x, \alpha).
\]
This function satisfies:
\begin{equation}
 \sigma(x, \alpha') = 0,\quad
\left|\frac{\partial \sigma}{\partial \alpha_1}\right| = \left|\frac{\frac{\partial \chi}{\partial \alpha_1}}{1 - \frac{\partial \chi}{\partial y}}\right|=0,\quad
\frac{\partial \sigma}{\partial \alpha_2}= \left|\frac{\frac{\partial \chi}{\partial \alpha_2}}{1 - \frac{\partial \chi}{\partial y}}\right|=\frac{1}{1-\lambda}.
\end{equation}

Consequently, for each $\alpha \in U(\alpha')$ and all $x \in \R$, we have the estimate
\[
|\sigma(x, \alpha)| \leq \frac{1}{1-\lambda}|\alpha_2 - \alpha'_2|.
\]


From \eqref{contrpo}, we obtain the non-expensive property:
\begin{align*}
|\pi_2 K_\alpha(x,y) - \pi_2 K_\alpha(x, \sigma(x, \alpha))|
&= |\chi(x,y,\alpha) - \chi(x, \sigma(x, \alpha), \alpha)| \\
&\leq \left|\frac{\partial \chi}{\partial y}\right| |y - \sigma(x, \alpha)| \\
&\leq |y - \sigma(x, \alpha)|.
\end{align*}

This shows that the set
\[
\Omega_1 := \{(x,y) \mid |y| \leq \frac{1}{1-\lambda}|\alpha_2 - \alpha'_2|\}
\]
is an attractor block for $K_\alpha$. Through conjugation, the corresponding set
\[
\Omega_2 := \{(x,y) \mid |y - \Psi_{\alpha'}(x)| \leq \frac{1}{1-\lambda}|\alpha_2 - \alpha'_2|\}
\]
is an attractor block for $F_\alpha$. Since the graph of $\Psi_\alpha$ is an attractor for $F_\alpha$, it must be contained in $\Omega_2$.

This completes the proof of Proposition~\ref{lipo}.
\end{proof}



Next, we consider the perturbed map defined by (\ref{Fp}). By Propositions \ref{hf1} and \ref{nhim}, if $\|\phi\|_{C^r}\leq \delta$, then the perturbed map $F^{\phi}_\alpha$ still admits a $C^r$-invariant graph ${\Gamma}_\alpha^\phi:=\{(x,{\Psi}_{\alpha}(x))\ |\ x\in \R\}$ and
\[F^{\phi}_\alpha(x,{\Psi}_{\alpha}(x))=(g_\al(x),{\Psi}_{\alpha}(g_\al(x))),\]
where $g_\al(x)=x+\al_1+\lambda{\Psi}_{\alpha}(x)+\phi(x)$. By Proposition \ref{lipo}, for each $x\in \R$, $g_\al(x)$ is Lipschitz continuous  with respect to $\al\in \R$.  Note that $g_\al\in D^r(\T)$. From the continuity of $\rho(f)$ with respect to $f\in D^r(\T)$, we obtain
\begin{Lemma}\label{contt}
$\rho(g_\al)$ is continuous with respect to $\al=(\al_1,\al_2)\in \R^2$.
\end{Lemma}


\subsection{Diversity of the dynamics on invariant graphs}
We now complete the proof of Theorem~\ref{M1} by establishing Item~(3).

Since $g_\alpha \in D^r(\T)$, there exists $\bar{g}_\alpha \in \mathrm{Diff}^r_+(\T)$ such that $g_\alpha$ is the lift of $\bar{g}_\alpha$ to the universal covering space. Let $\mu_\alpha$ be an ergodic invariant probability measure for $\bar{g}_\alpha$ on $\T$. By \cite[Proposition 2.3]{H11}, the rotation number satisfies
\begin{equation}\label{rotation}
\rho(g_\alpha) = \alpha_1 + \lambda \int_{\T} \Psi_\alpha \, d\mu_\alpha+\int_{\T} \phi \, d\mu_\alpha.
\end{equation}

From Proposition~\ref{nhim}, we have the uniform estimate
\[
\left\|\Psi_\alpha(x) - \frac{\alpha_2}{1-\lambda}\right\|_{C^1} \leq \varepsilon.
\]
This allows us to rewrite \eqref{rotation} as
\begin{equation}\label{rotation-rewritten}
\rho(g_\alpha) = \alpha_1 + \frac{\lambda\alpha_2}{1-\lambda} + \lambda \int_{\T} \left(\Psi_\alpha(x) - \frac{\alpha_2}{1-\lambda}\right) d\mu_\alpha+\int_{\T} \phi \, d\mu_\alpha.
\end{equation}

For fixed $\alpha_2 \in \R$, define the function $h_{\alpha_2}:\R\to\R$ by
\[
h_{\alpha_2}(\alpha_1) := \alpha_1 + \frac{\lambda\alpha_2}{1-\lambda} + \lambda \int_{\T} \left(\Psi_{(\alpha_1,\alpha_2)}(x) - \frac{\alpha_2}{1-\lambda}\right) d\mu_\alpha+\int_{\T} \phi \, d\mu_\alpha.
\]

By Lemma~\ref{contt}, the function $h_{\alpha_2}$ is continuous in $\alpha_1 \in \R$. Moreover, we observe that:
\begin{itemize}
\item For any fixed $\alpha_2$, $h_{\alpha_2}$ is surjective since
\[
\lim_{\alpha_1 \to \pm\infty} h_{\alpha_2}(\alpha_1) = \pm\infty,
\]
and the integral term remains bounded by $\lambda\varepsilon$.

\item Consequently, for any $\omega \in \R$, there exists $\bar{\alpha}_1 \in \R$ such that $\rho(g_{(\bar{\alpha}_1,\alpha_2)}) = \omega$.

\item Similarly, for fixed $\alpha_1$ and any $\omega \in \R$, there exists $\bar{\alpha}_2 \in \R$ satisfying $\rho(g_{(\alpha_1,\bar{\alpha}_2)}) = \omega$.
\end{itemize}

This completes the proof of Theorem~\ref{M1}.



\subsection{Generic frequency changes for a fixed system}
We omit the superscript $\tilde{A}$ in the notation $A$ for simplicity and prove Theorem \ref{Mt3} here.

Inspired by the Arnold family \cite{Ar}, we consider, for a given $\alpha_1 \in \R$, the following map:
\begin{equation}\label{g22}
g_{n,\alpha_1}(x) := x + \alpha_1 + \frac{1}{n} \sin(2\pi x).
\end{equation}
For sufficiently large $n$ (say $n\geq n_0$), we have $g_{n,\alpha_1} \in D^\omega(\T)$. Let $g_{n,\alpha_1}^{-1}(x) := x - \alpha_1 - \xi_n(x)$ denote its inverse. Due to the symmetry of the graphs of $g_{n,\alpha_1}$ and $g^{-1}_{n,\alpha_1}$, we obtain
\[
\int_{\T} \xi_n(x) \, dx = \int_{\T} \frac{1}{n} \sin(2\pi x) \, dx = 0.
\]

%Define the function
%\[
%\psi_n(x) := \frac{1}{\lambda n} \sin(2\pi x) - \xi_n(x).
%\]
By Proposition~\ref{hf1}, the map $F_{\alpha_1}^{\psi_n}$ defined in \eqref{Fp} admits a $C^0$-invariant graph
\[
\Gamma_n := \{(x, \Psi_{n,\alpha_1}(x)) \mid x \in \R\},
\]
if and only if
\begin{equation}\label{g+g}
\frac{1}{1+\lambda}g_{n,\alpha_1}(x) + \frac{\lambda}{1+\lambda}g_{n,\alpha_1}^{-1}(x) = x + \frac{1}{1+\lambda}\left((1-\lambda)\alpha_1  + \psi_n(x)\right).
\end{equation}
We construct
\[\Psi_{n,\alpha_1}(x):=\frac{1}{n} \sin(2\pi g_{n,\al_1}^{-1}(x)),\quad \psi_n(x):=\Psi_{n,\alpha_1}(g_{n,\al_1}(x))-\lambda\Psi_{n,\alpha_1}(x).\]
Then (\ref{g+g}) holds. More precisely, we have
\[
\psi_n(x) = \frac{1}{n} \sin(2\pi x) - \frac{\lambda}{n} \sin(2\pi g_{n,\alpha_1}^{-1}(x)).
\]

We observe the following bounds:
\[
\frac{1}{2} \leq \|Dg_{n,\alpha_1}\|_{C^0} \leq 2, \quad \text{and} \quad \|Dg_{n,\alpha_1}^{-1}\|_{C^0} = \frac{1}{\|Dg_{n,\alpha_1}\|_{C^0}} \leq 2.
\]
By the  Fa\`{a} di Bruno formula (see \cite[Corollary 2.4]{H11} for instance), for each integer $r \geq 1$, there exists a constant $B_r$ depending only on $r$ such that
\begin{itemize}
\item if $r=1$, $\|D^{r} g_{n,\alpha_1}^{-1}\|_{C^0}\in [1-\frac{B_1}{n}, 1+\frac{B_1}{n}]$;
\item if $r\geq 2$, $\|D^{r} g_{n,\alpha_1}^{-1}\|_{C^0}\leq \frac{B_r}{n}$.
\end{itemize}

A direct calculation yields the existence of a constant $C_r$, depending only on $r$, satisfying
\[
\|\psi_n\|_{C^r} \leq \frac{C_r}{n}.
\]

To complete the argument, we construct a $G_\delta$ subset $\mathcal{O} \subset \R$ such that for every $\alpha_1 \in \mathcal{O}$, the circle map induced by the restriction of $F^{\psi_n}_{\alpha_1}$ to $\Gamma_n$ has rotation number different from $\alpha_1$.

According to \cite[Lemma 4.2]{de}, the set
\[
E_n := \{\alpha_1 \in \R \mid \rho(g_{n,\alpha_1}) \text{ is irrational}\}
\]
is nowhere dense in $\R$. Let $O_n := \R \setminus \overline{E}_n$ denote its complement, which is open and dense in $\R$. We can therefore express $O_n$ as a countable union of open intervals:
\[
O_n = \bigcup_{k=1}^{\infty} I_k.
\]
Letting
\[
{O}:= \bigcap_{n=n_0}^{\infty} {O}_n.
\]
Then the set $O$ is a $G_\delta$ set.


Applying \cite[Lemma 4.2]{de} again, we find that for each $\alpha_1 \in I_k$, the rotation number $\rho(g_{n,\alpha_1})$ equals some fixed rational number $r_k$ only depending on $k$. Note that the set of irrational numbers $\R\backslash\mathbb{Q}$ is also a $G_\delta$ set. We then define
\[\mathcal{O} := O\cap\R\backslash\mathbb{Q}.\]
It follows that $\mathcal{O}$ is still a $G_\delta$ set. For each $\al_1\in \mathcal{O}$, $\al_1$ is irrational. However,  the rotation number $\rho(g_{n,\alpha_1})$ is rational for all $n\geq n_0$. Therefore, the set $\mathcal{O}$ is the desired dense $G_\delta$ subset of $\R$ for which the theorem holds.






 \vspace{2em}

	\section{\sc The threshold for the existence of invariant graphs}\label{S4}
	

\subsection{Persistence under Lipschitz perturbations}
We give a proof of Theorem \ref{M2} here.
\subsubsection{Construction of the graph transform}
Fix $\la\in (0,1)$. Let $I := [-(1-\lambda), 1-\lambda]$ and denote by $\Lip(\T, I)$ the space of 1-periodic Lipschitz maps $\psi: \R \to I$ with Lipschitz constant $\mathcal{L}(\psi)$. For $k > 0$, define
\[
\Lip_k := \{\psi \in \Lip(\T, I) \mid \mathcal{L}(\psi) \leq k\}.
\]

Consider the map $F_{\al_1}^{\phi^\lambda}(x,y) = (X(x,y), Y(x,y))$ where
\begin{align*}
X(x,y) &:= x + \alpha_1 + \lambda y + \phi^\lambda(x), \\
Y(x,y) &:= \lambda y + \phi^\lambda(x).
\end{align*}
For any $\psi \in \Lip_k$, we have the compositions
\begin{align*}
X \circ (\mathrm{Id}, \psi)(x) &= x + \alpha_1 + \lambda \psi(x) + \phi^\lambda(x), \\
Y \circ (\mathrm{Id}, \psi)(x) &= \lambda \psi(x) + \phi^\lambda(x).
\end{align*}
Let $A_\lambda := \|\phi^\lambda\|_{\mathrm{Lip}}$. It follows that $\|\phi^\la\|_{C^0}\leq A_\lambda$.

\begin{Lemma}\label{invertib}
Define the constant
\[
K_1 := \frac{2}{\sqrt{\lambda}} - 1.
\]
If $k < K_1$, then for each $\psi \in \Lip_k$, the map $X \circ (\mathrm{Id}, \psi)$ is invertible with Lipschitz inverse satisfying
\[
\mathcal{L}\left([X \circ (\mathrm{Id}, \psi)]^{-1}\right) \leq \frac{1}{1 - \lambda k - A_\lambda}.
\]
\end{Lemma}

\begin{proof}
Define $u(x) := X \circ (\mathrm{Id}, \psi)(x) - x$. Then $u$ is Lipschitz with
\[
|u(x_1) - u(x_2)| \leq (\lambda k + A_\lambda)|x_1 - x_2|.
\]
The map $X \circ (\mathrm{Id}, \psi)$ is invertible if $\mathcal{L}(u) < 1$, in which case
\[
\mathcal{L}\left([X \circ (\mathrm{Id}, \psi)]^{-1}\right) \leq \frac{1}{1 - \mathcal{L}(u)}.
\]
The condition $k < K_1$ implies $\lambda k + (1 - \sqrt{\lambda})^2 < 1$, and consequently
\[
\mathcal{L}(u) \leq \lambda k + A_\lambda \leq \lambda k + (1 - \sqrt{\lambda})^2 < 1.
\]
\end{proof}

When $[X \circ (\mathrm{Id}, \psi)]^{-1}$ exists, we define the graph transform $\mathcal{T}: \Lip_k \to \Lip_k$ by
\[
\mathcal{T}\psi : x \mapsto Y \circ (\mathrm{Id}, \psi) \circ [X \circ (\mathrm{Id}, \psi)]^{-1}(x).
\]
For the graph $\tilde{\Gamma}(\psi) := \{(x, \psi(x)) \mid x \in \R\}$, we have the invariance property
\[
\tilde{\Gamma}(\mathcal{T}\psi) = F(\tilde{\Gamma}(\psi)).
\]

\begin{Lemma}\label{welld}
Define the constant
\[
K_2 := \frac{1}{\sqrt{\lambda}} - 1.
\]
The graph transform $\mathcal{T}: \Lip_{K_2} \to \Lip_{K_2}$ is well-defined, i.e., $\mathcal{T}\psi \in \Lip_{K_2}$ for all $\psi \in \Lip_{K_2}$.
\end{Lemma}

\begin{proof}
For $\psi \in \Lip_{K_2}$, we first establish the uniform bound:
\[
\|\mathcal{T}\psi\|_{C^0} \leq \lambda \|\psi\|_{C^0} + A_\lambda \leq \lambda(1-\lambda) + (1 - \sqrt{\lambda})^2 < 1-\lambda.
\]

The Lipschitz estimate follows from:
\begin{align*}
|\mathcal{T}\psi(x_1) - \mathcal{T}\psi(x_2)|
&\leq \frac{\lambda \mathcal{L}(\psi) + A_\lambda}{1 - \lambda \mathcal{L}(\psi) - A_\lambda}|x_1 - x_2| \\
&\leq \frac{\lambda K_2 + (1 - \sqrt{\lambda})^2}{1 - \lambda K_2 - (1 - \sqrt{\lambda})^2}|x_1 - x_2| \\
&= K_2 |x_1 - x_2|.
\end{align*}
Thus $\mathcal{T}\psi \in \Lip_{K_2}$.
\end{proof}




\subsubsection{Contraction mapping}
Note that $\Lip_k$ is a closed subspace of the Banach space $C^0(\T,I)$ equipped with the $C^0$-metric, and hence is complete. To complete the proof of Theorem~\ref{M2}, it remains to show that the graph transform is a contraction mapping. Assuming the invertibility of $X \circ (\Id,\psi)$ and the well-definedness of $\mathcal{T}$, we establish this through the following lemma.

\begin{Lemma}\label{ctrmap}
Define the constant
\[
K_3 := \frac{1}{\lambda} - 1.
\]
If $k < K_3$, then the graph transform $\mathcal{T}: \Lip_k \to \Lip_k$ is a contraction. Specifically, for any $\psi_1, \psi_2 \in \Lip_k$,
\[
\|\mathcal{T}\psi_1 - \mathcal{T}\psi_2\|_{C^0} \leq l \|\psi_1 - \psi_2\|_{C^0},
\]
where $0 < l < 1$.
\end{Lemma}

\begin{proof}
Fix $z \in \T$ and $\psi_1, \psi_2 \in \Lip_k$. Let $(x,y)$ be the point on the graph of $\psi_1$ determined by
\[
x := [X \circ (\Id, \psi_1)]^{-1}(z), \quad y := \psi_1(x).
\]
By definition of the graph transform, we have:
\begin{align*}
\mathcal{T}\psi_1(z) &= Y(\Id, \psi_1)(x), \\
\mathcal{T}\psi_2(z) &= \mathcal{T}\psi_2 \circ X(x, \psi_1(x)) = Y(\Id, \psi_2)(x'),
\end{align*}
where $x' := [X \circ (\Id, \psi_2)]^{-1} \circ X(x, \psi_1(x))$.

The difference can be estimated as:
\begin{align*}
|\mathcal{T}\psi_1(z) - \mathcal{T}\psi_2(z)|
&\leq |Y(\Id, \psi_1)(x) - Y(\Id, \psi_2)(x)| \\
&\quad + |Y(\Id, \psi_2)(x) - \mathcal{T}\psi_2 \circ X(\Id, \psi_1)(x)| \\
&\leq \lambda |\psi_1(x) - \psi_2(x)| + k\lambda |\psi_1(x) - \psi_2(x)| \\
&= (\lambda + k\lambda) |\psi_1(x) - \psi_2(x)|.
\end{align*}

This yields the uniform estimate:
\[
\|\mathcal{T}\psi_1 - \mathcal{T}\psi_2\|_{C^0} \leq l \|\psi_1 - \psi_2\|_{C^0},
\]
where $l := \lambda(1 + k)$. When $k < K_3 = \frac{1}{\lambda} - 1$, we have $l < 1$, proving that $\mathcal{T}$ is indeed a contraction.
\end{proof}











Note that for each $\lambda\in (0,1)$, we have
\[K_2<\min\{K_1,K_3\}.\]
It follows that the map $\mathcal{T} \colon \mathrm{Lip}_{K_2} \to \mathrm{Lip}_{K_2}$ admits a unique fixed point $\Psi$. Moreover, the graph
\begin{equation}\label{tGp}
\tilde{\Gamma}(\Psi) := \bigl\{ \bigl(x, \Psi(x)\bigr) \bigm| x \in \mathbb{R} \bigr\}
\end{equation}
is the unique Lipschitz invariant graph for $F_{\al_1}^{\phi^\lambda}$.


 \vspace{2em}



\subsection{Persistence under $C^1$ perturbations}
We prove Theorem \ref{M22} here. The proof is inspired by \cite[Section 4]{HPS} and \cite[Proof of Theorem 3.1]{BB}).
\subsubsection{The cone condition}
Denote
\[\beta:=\frac{1}{\sqrt{\lambda}}-1.\]
 We denote $F$ for $F_{\al_1}^{\phi^\lambda}$, and we use $\Gamma$ instead of $\tilde{\Gamma}$ to denote the invariant graph by $F$  for simplicity.
For $z\in\R^2$, we consider the cone  as follows:
\[\mathcal{C}_\beta(z):=\{v=(v_1,v_2)\in\R^2\ |\ |v_2|\leq \beta |v_1|\},\]
where we identified the Euclidean space $\R^2$ and its tangent space $\mathrm{T}_z\R^2$. The inner product in $\R^2$ is given by the standard one. Recall the projection $\pi_1:\R^2\to\R$ via
\[\pi_1:(x,\Psi(x))\mapsto x,\]
which means $\pi_1^{-1}(x)\in \Gamma$ for each $x\in\R$. Moreover, we denote the cone along $\Gamma$ by $\mathcal{C}_\beta(x)$ instead of  $\mathcal{C}_\beta(\pi_1^{-1}(x))$ for simplicity. Since the invariant graph $\Gamma$ is {\it a priori} only Lipschitz, we also need to consider the tangent cone along $\Gamma=\{(x,\Psi(x))\ |\ x\in\R\}$.
\[TC_\Gamma(x):=\left\{v\in\R^2\ \left|\ v=\mu\lim_{n\to+\infty}\frac{(x_n,\Psi(x_n))^T-(x,\Psi(x))^T}{\|(x_n,\Psi(x_n))^T-(x,\Psi(x))^T\|}\right.,\quad \forall \mu\in\R\right\},\]
where  $a^T$ denotes the transpose of $a$. Let us recall that there exists $g\in D^0(\T)$ such that
\[F(x,\Psi(x))=(g(x),\Psi(g(x))).\]
A direct calculation implies that for each $v\in TC_\Gamma(x)$, there exists $\mu\in\R$ such that
\[DF(\pi_1^{-1}(x))v=\mu\lim_{n\to+\infty}\frac{(g(x_n),\Psi(g(x_n)))^T-(g(x),\Psi(g(x)))^T}{\|(g(x_n),\Psi(g(x_n)))^T-(g(x),\Psi(g(x)))^T\|}.\]
It follows that
\begin{equation}\label{cooncd}
DF(\pi_1^{-1}(x))\cdot\left(TC_\Gamma(x)\right):=\left\{DF(\pi_1^{-1}(x))v\ |\ v\in TC_\Gamma(x)\right\}\subseteq TC_\Gamma(g(x)),
\end{equation}

Following \cite[Proof of Theorem 3.1]{BB}), we need to verify the following two conditions:
\begin{itemize}
\item [(1)] $TC_\Gamma(x)\subseteq \mathcal{C}_\beta(x)$ for each $x\in \R$;
\item [(2)] $DF(\pi_1^{-1}(x))\cdot(\mathcal{C}_\beta(x))\subseteq \mathrm{interior}\ \mathcal{C}_\beta(x)\cup \{(0,0)\}$  for each $x\in \R$.
\end{itemize}
By Lemma \ref{welld} and (\ref{tGp}), we know that $\mathcal{L}(\Psi)\leq \frac{1}{\sqrt{\lambda}}-1$, which means Item (1) holds.
For Item (2), it follows directly from the definition of $F$ that for any vector $v = (v_1, v_2) \in \mathcal{C}_\beta(x)$, we have
\[
DF(\pi_1^{-1}(x))v =
\begin{pmatrix}
1+\phi'(x) & \lambda \\
\phi'(x) & \lambda
\end{pmatrix}
\begin{pmatrix}
v_1 \\
v_2
\end{pmatrix}
=
\begin{pmatrix}
v_1 + \phi'(x)v_1 + \lambda v_2 \\
\phi'(x)v_1 + \lambda v_2
\end{pmatrix}.
\]
Note that $\|\phi\|_{C^1} < (1 - \sqrt{\lambda})^2$ and $|v_2 / v_1| \leq \beta = \frac{1}{\sqrt{\lambda}} - 1$. It is straightforward to verify that
\begin{equation}\label{test2}
\frac{|v_1 + \phi'(x)v_1 + \lambda v_2|}{|\phi'(x)v_1 + \lambda v_2|}
= \frac{|1 + \phi'(x) + \lambda \frac{v_2}{v_1}|}{|\phi'(x) + \lambda \frac{v_2}{v_1}|}
< \beta,
\end{equation}
which confirms the validity of Item (2).


\subsubsection{Differentiability}
Let us recall a classical result regarding the geometrical criterion on the differentiability of a Lipschitz submanifold (see \cite[Lemma 4.2]{BB} for instance).
\begin{Lemma}\label{difli}
Let \( Z \) be a Lipschitz submanifold of dimension \( n \). If for every \( z \in Z \), the tangent cone \( TC_Z(z) \) is contained in an \( n \)-dimensional space \( L(z) \), then \( Z \) is a differentiable submanifold with \( T_zZ = L(z) \). If moreover \( z \mapsto L(z) \) is continuous, then \( Z \) is of class \( C^1 \).
\end{Lemma}

In view of  Item (2) in last subsection, we have (see \cite[Proof of Theorem 1.2]{Ne}) for each $x\in\R$, the set
\[L(x):=\cap_{k\geq 0}DF^{k}(\pi_1^{-1}(g^{-k}(x)))\cdot (\mathcal{C}_\beta(g^{-k}(x)))\subseteq \mathcal{C}_\beta(x)\]
is a 1-dimensional subspace of $T_{\pi_1^{-1}(x)}\R^2$.

By (\ref{cooncd}), there holds for each $x\in\R$,
\[TC_{\Gamma}(x)=\cap_{k\geq 0}DF^{k}(\pi_1^{-1}(g^{-k}(x)))\cdot (TC_{\Gamma}(g^{-k}(x))).\]
By construction, we have
\[\cup_{x\in\R}TC_{\Gamma}(x)\subseteq \cup_{x\in\R}\mathcal{C}_\beta(x),\]
which yields from the definition of $L(x)$ that
\[\cup_{x\in\R}TC_{\Gamma}(x)\subseteq \cup_{x\in\R}L(x).\]
By Lemma \ref{difli}, the function $\Psi$ is differentiable, with $T_{\pi_1^{-1}(x)}\Gamma=L(x)$ for each $x\in \R$.





\subsubsection{Continuous differentiability}
We prove that $\Psi$ is continuously differentiable. Given a convergent sequence \( x_n \to x \) in $\mathbb{R}$, we need to show that
\[
L(x_n) := T_{\pi_1^{-1}(x_n)}\Gamma
\]
converges to
\[
L(x) := T_{\pi_1^{-1}(x)}\Gamma
\]
in the Hausdorff topology. By compactness of the Grassmannian, it suffices to show that \( L(x) \) is the unique accumulation point of the sequence \( L(x_n) \).

Assume, for the sake of contradiction, that there exists an accumulation point \( L'(x) \neq L(x) \) of the sequence \( L(x_n) \). Note that for each \( n \), we have
\[
L(x_n) \subseteq \mathcal{C}_\beta(x_n),
\]
where \( \mathcal{C}_\beta(x) \) is continuous with respect to \( x \in \mathbb{R} \) in the Hausdorff topology. Taking limits, it follows that
\[
L'(x) \subseteq \mathcal{C}_\beta(x).
\]

For each \( k \geq 0 \), the continuity of \( DF^{-k} \) with respect to \( x \in \mathbb{R} \) implies that
\[
DF^{-k}(\pi_1^{-1}(x)) \cdot L'(x)
\]
is an accumulation point of the sequence
\[
DF^{-k}(\pi_1^{-1}(x_n)) \cdot L(x_n) = T_{\pi_1^{-1}(g^{-k}(x_n))}\Gamma.
\]
Again, by continuity of \( \mathcal{C}_\beta \), we have:
\[
DF^{-k}(\pi_1^{-1}(x)) \cdot L'(x) \subseteq \mathcal{C}_\beta(g^{-k}(x)).
\]
By the definition of \( L(x) \), it follows that:
\[
L'(x) \subseteq \bigcap_{k \geq 0} DF^k(\pi_1^{-1}(g^{-k}(x))) \cdot \mathcal{C}_\beta(g^{-k}(x)) = L(x).
\]
Since \( L'(x) \) and \( L(x) \) have the same dimension, we must have \( L'(x) = L(x) \), a contradiction. Therefore, \( L(x) \) is the unique accumulation point of \( L(x_n) \), and we conclude that \( x \mapsto L(x) \) is continuous in the Hausdorff topology.

Thus, $\Psi$ is continuously differentiable.

\begin{Remark}\label{needb}
In the proof of Theorem \ref{M222} below, we require a slight generalization of Theorem \ref{M22}. Given parameters $\lambda_1, \lambda_2 \in (0,1)$, consider the map
\begin{equation}\label{Fpwp}
F^\phi(x,y) = \big(x + \alpha_1 + \lambda_1 y + \phi_1(x),\; \lambda_2 y + \phi_2(x)\big).
\end{equation}
Let $\tilde{\lambda} := \max\{\lambda_1, \lambda_2\}$. Then for every $C^1$ perturbation $\phi_1, \phi_2$ satisfying
\[
\max\left\{\|\phi_1\|_{C^1},\|\phi_2\|_{C^1}\right\} < (1 - \sqrt{\tilde{\lambda}})^2,
\]
the map $F^\phi$ admits a unique $C^1$ invariant graph.

To prove this, it suffices to verify that if we set $K_2 := \frac{1}{\sqrt{\tilde{\lambda}}} - 1$, then the following inequality holds:
\[
\frac{\lambda_2 K_2 + (1 - \sqrt{\tilde{\lambda}})^2}{1 - \lambda_1 K_2 - (1 - \sqrt{\tilde{\lambda}})^2} \leq K_2.
\]
This ensures that the cone condition is preserved, allowing us to apply the same argument as in the proof of Theorem \ref{M22}, including the verification of inequality~\eqref{test2}.

Indeed, we consider two cases:

\medskip
\noindent\textbf{Case 1:} $\lambda_1 \leq \lambda_2$. Then $\tilde{\lambda} = \lambda_2$, and we compute
\[
\frac{\lambda_2 K_2 + (1 - \sqrt{\lambda_2})^2}{1 - \lambda_1 K_2 - (1 - \sqrt{\lambda_2})^2}
\leq \frac{\lambda_2 K_2 + (1 - \sqrt{\lambda_2})^2}{1 - \lambda_2 K_2 - (1 - \sqrt{\lambda_2})^2}
\leq K_2.
\]

\noindent\textbf{Case 2:} $\lambda_1 > \lambda_2$. Then $\tilde{\lambda} = \lambda_1$, and we compute
\[
\frac{\lambda_2 K_2 + (1 - \sqrt{\lambda_1})^2}{1 - \lambda_1 K_2 - (1 - \sqrt{\lambda_1})^2}
\leq \frac{\lambda_1 K_2 + (1 - \sqrt{\lambda_1})^2}{1 - \lambda_1 K_2 - (1 - \sqrt{\lambda_1})^2}
\leq K_2.
\]

In both cases, the inequality holds, completing the proof of the generalization.
\end{Remark}


 \vspace{2em}

\section{\sc The normal form and sharpness of $\gm$}
The proof of Theorem \ref{M222} is inspired by \cite{Ma}.
Fix $\alpha_2 = 0$ and assume that $\alpha_1$ is a Diophantine number. We rewrite the model defined by \eqref{Fp} as
\begin{equation}\label{Fpw}
F^\phi_{\lambda}(x, y) = \big(x + \alpha_1 + \lambda y + \phi(x),\; \lambda y + \phi(x)\big),
\end{equation}
where $\phi$ is a $1$-periodic function satisfying the normalization condition
\[
\int_\T \phi(x)\,dx = 0.
\]
Let $C^\omega(\mathbb{T})$ denote the set of real-analytic functions on the circle $\mathbb{T}$. Given any $\phi \in C^\omega(\mathbb{T})$, there exists $s := s(\phi) > 0$ such that $\phi$ admits a holomorphic extension to the complex strip
\[
\mathbb{T}_s := \left\{ \theta \in \mathbb{C}/\mathbb{Z} \ \middle| \ |\mathrm{Im}\,\theta| \leq s \right\}.
\]
We say that $\phi:\mathbb{T}_s \to \mathbb{C}$ is well-defined if there exists a unique holomorphic extension $\bar{\phi}$ to $\mathbb{T}_s$ with finite Banach norm
\begin{equation}\label{renorm}
\|\phi\|_s := \sup_{\theta \in \mathbb{T}_s} |\bar{\phi}(\theta)|.
\end{equation}
Hereafter, we will use the notation $\|\cdot\|_{C^\omega}$ instead of $\|\cdot\|_s$ whenever the analytic radius does not need to be emphasized.


\subsection{R\"{u}ssmann's normal form}

By R\"{u}ssmann's normal form \cite{R1} (see also \cite[Theorem 5.4]{Mas1} for a higher-dimensional generalization), we have the following result:

\begin{Proposition}\label{prrusss}
There exists $\varepsilon_0 > 0$ such that for each $\varepsilon \in (0, \varepsilon_0]$, if $\|\phi\|_{C^\omega} \leq \varepsilon$, then there exists a function $\Psi_\lambda \in C^\omega(\mathbb{T})$ and a map $h_\lambda \in D^\omega(\mathbb{T})$ with $h_\lambda(0) = 0$ such that
\begin{equation}\label{fhold}
F^\phi_\lambda(x, y) = \big(h_\lambda \circ R_{\alpha_1} \circ h_\lambda^{-1}(x),\; \nu_\lambda + \Psi_\lambda(h_\lambda \circ R_{\alpha_1} \circ h_\lambda^{-1}(x))\big).
\end{equation}
Moreover,
\[
\|\Psi_\lambda\|_{C^\omega} = O(\varepsilon), \quad \|h_\lambda - \mathrm{Id}\|_{C^\omega} = O(\varepsilon), \quad \nu_\lambda = O(\varepsilon).
\]
\end{Proposition}

\begin{Remark}
We refer to \cite[Theorem A.1]{Mas1} for a more quantitative estimate of the $O(\varepsilon)$ terms. In particular, from the proof of \cite[Theorem 5.1]{Mas1}, one can see that for $\lambda$ bounded away from zero (e.g., $\lambda \in [\frac{1}{2}, 1]$), the constants involved in the $O(\varepsilon)$ terms can be taken independently of $\lambda$.
\end{Remark}

Proposition \ref{prrusss} leads to the following corollary:

\begin{Lemma}
Let $\gamma \in (1, 2)$. Then there exists $\lambda_0 \in (0, 1)$ such that for every $\lambda \in [\lambda_0, 1)$, if
\[
\|\phi\|_{C^\omega} \leq (1 - \lambda)^\gamma,
\]
then there exists a function $\Psi_\lambda \in C^\omega(\mathbb{T})$ and a map $h_\lambda \in D^\omega(\mathbb{T})$ with $h_\lambda(0) = 0$ such that \eqref{fhold} holds, and
\[
\|\Psi_\lambda\|_{C^\omega} = O((1 - \lambda)^\gamma), \quad \|h_\lambda - \mathrm{Id}\|_{C^\omega} = O((1 - \lambda)^\gamma), \quad \nu_\lambda = O((1 - \lambda)^\gamma).
\]
\end{Lemma}

In the following, we shall denote $\varepsilon := 1 - \lambda$ and $\alpha := \alpha_1$, and omit the subscript $\lambda$ for simplicity.


\subsection{Changes of coordinates}

The proof of Theorem \ref{M222} will be completed through a sequence of coordinate transformations. We begin by introducing the change of variables
\[
H: (x, y) \mapsto (\xi, \eta) \quad \text{via} \quad
\left\{
  \begin{array}{l}
    \xi = h^{-1}(x), \\
    \eta = y - \Psi(x),
  \end{array}
\right.
\]
where $h$ and $\Psi$ are as given in Proposition \ref{prrusss}. Then the transformed map $\bar{F}^\phi := H \circ F^{\phi} \circ H^{-1}$ takes the form
\begin{align*}
\bar{F}^\phi(\xi, \eta)
&= \left(h^{-1}(h(\xi + \alpha) + \lambda \eta),\; \nu + \Psi(h(\xi + \alpha)) + \lambda \eta - \Psi(h(\xi + \alpha) + \lambda \eta)\right).
\end{align*}

Since $h \in D^\omega(\mathbb{T})$, we have $h^{-1} \in D^\omega(\mathbb{T})$ and
\[
\|h^{-1} - \mathrm{Id}\|_{C^\omega} = O(\varepsilon).
\]
Thus, there exists $\bar{s} > 0$, depending on the analyticity radii of $\phi$ and $\Psi$, such that for all $\xi \in \mathbb{T}$ and $|\lambda \eta| \leq \bar{s}$, we have the expansion:
\begin{equation}
\bar{F}^\phi(\xi, \eta) = \left(\xi + \alpha + \sum_{i=1}^\infty A_i(\xi)(\lambda \eta)^i,\; \nu + \sum_{i=1}^\infty B_i(\xi)(\lambda \eta)^i\right),
\end{equation}
where
\[
A_i(\xi) = \frac{1}{i!} D^i h^{-1}(h(\xi + \alpha)), \qquad
B_i(\xi) = -\frac{1}{i!} D^i \Psi(h(\xi + \alpha)) \quad \text{for } i \geq 2,
\]
and
\[
A_1(\xi) = 1 + O(\varepsilon^\gamma), \quad B_1(\xi) = 1 - D\Psi(h(\xi + \alpha)) = 1 + O(\varepsilon^\gamma).
\]
In particular,
\[
A_i(\xi) = O(\varepsilon^\gamma), \quad B_i(\xi) = O(\varepsilon^\gamma) \quad \text{for all } i \geq 2.
\]

\begin{Lemma}\label{ketra}
Let $\alpha$ be a Diophantine number. Then there exists $\tilde{s} > 0$ and a real-analytic change of coordinates $T: \mathbb{R}^2 \to \mathbb{R}^2$ given by $(\xi, \eta) \mapsto (X, Y)$ such that for each $X \in \mathbb{T}$ and $|\lambda Y| \leq \tilde{s}$, the transformed map
\[
\tilde{F}^\phi := T \circ \bar{F}^\phi \circ T^{-1} : (X, Y) \mapsto (\tilde{X}, \tilde{Y})
\]
satisfies
\[
\left\{
\begin{array}{l}
\tilde{X} = X + \alpha + \bar{\alpha}_1 \cdot (\lambda Y) + \bar{\alpha}_2 \cdot (\lambda Y)^2 + O(\varepsilon^\gamma |\lambda Y|^3) + O(\varepsilon^\gamma |\nu|), \\
\tilde{Y} = \nu + \bar{\beta}_1 \cdot (\lambda Y) + \bar{\beta}_2 \cdot (\lambda Y)^2 + O(\varepsilon^\gamma |\lambda Y|^3) + O(\varepsilon^\gamma |\nu|).
\end{array}
\right.
\]
\end{Lemma}

The following technical result is taken from \cite[Lemma B.1]{Ma} and will be used in the construction of the coordinate change $T$:

\begin{Lemma}\label{jsl}
Let $\alpha$ be a $(D, \tau)$-Diophantine number (as defined in \eqref{diop}). Let $\mathcal{A}(\mathbb{T}_s)$ denote the space of holomorphic functions $g: \mathbb{T}_s \to \mathbb{C}$. Given constants $a, b \in \mathbb{R}$ with $b \neq 0$, and given $0 < \sigma < s$, for each $g \in \mathcal{A}(\mathbb{T}_s)$, there exists a unique $f \in \mathcal{A}(\mathbb{T}_{s - \sigma})$ with zero average, i.e., $\int_\mathbb{T} f(\theta)\, d\theta = 0$, and a unique constant $\mu \in \mathbb{R}$ such that
\[
\mu + a f(\theta + \alpha) - b f(\theta) = g(\theta), \quad \text{with} \quad \mu = \int_\mathbb{T} g(\theta)\, d\theta,
\]
and the solution satisfies the estimate
\[
\|f\|_{s-\sigma} \leq \frac{C}{D} \sigma^{-(\tau + 3)} \|g\|_s,
\]
where $C$ is a constant depending only on $\tau$.
\end{Lemma}




\subsubsection{Step 1: Elimination of the non-constant linear term}

According to Lemma~\ref{jsl}, by taking $a = b = 1$, the cohomological equation
\[
\int_{\mathbb{T}} \ln B_1(\xi)\, d\xi = \ln B_1(\xi) + U_1(\xi) - \ln U_1(\xi + \alpha)
\]
admits a unique solution $U_1 \in C^\omega(\mathbb{T})$. Since $B_1(\xi) = 1 + O(\varepsilon^\gamma)$, we deduce that
\[
U_1(\xi) = 1 + O(\varepsilon^\gamma), \quad \text{for all } \xi \in \mathbb{T}.
\]

We now define a coordinate transformation $T_1: (\xi, \eta) \mapsto (\xi, \eta_1)$ by
\[
\eta_1 = \frac{\eta}{U_1(\xi)},
\]
so that $T_1^{-1}(\xi, \eta_1) = (\xi, U_1(\xi)\eta_1)$. Let $\bar{\beta}_1 := \int_\mathbb{T} B_1(\xi)\, d\xi$. Then, for $|\lambda \eta_1| \leq \bar{s}$, the transformed map is given by
\[
\bar{F}^\phi_1 := T_1 \circ \bar{F}^\phi \circ T_1^{-1} : (\xi, \eta_1) \mapsto (\bar{\xi}, \bar{\eta}_1),
\]
where
\[
\left\{
\begin{array}{ll}
\bar{\xi} = \xi + \alpha + \alpha_1(\xi)(\lambda \eta_1) + \alpha_2(\xi)(\lambda \eta_1)^2 + O(\varepsilon^\gamma |\lambda \eta_1|^3), \\
\bar{\eta}_1 = \nu + \bar{\beta}_1 \cdot (\lambda \eta_1) + \beta_2(\xi)(\lambda \eta_1)^2 + O(\varepsilon^\gamma |\lambda \eta_1|^3) + O(\varepsilon^\gamma |\nu| |\lambda \eta_1|) + O(\varepsilon^\gamma |\nu|).
\end{array}
\right.
\]
Without loss of generality, we may assume $\bar{s} \leq 1$, and simplify the above to
\[
\bar{\eta}_1 = \nu + \bar{\beta}_1 \cdot (\lambda \eta_1) + \beta_2(\xi)(\lambda \eta_1)^2 + O(\varepsilon^\gamma |\lambda \eta_1|^3) + O(\varepsilon^\gamma |\nu|).
\]

Next, define $\bar{\alpha}_1 := \int_\mathbb{T} \alpha_1(\xi)\, d\xi$, and consider the change of coordinates $S_1: (\xi, \eta_1) \mapsto (\xi_1, \eta_1)$ given by
\[
\xi_1 = \xi + V_1(\xi) \eta_1,
\]
where $V_1(\xi)$ is the unique solution to the cohomological equation
\[
V_1(\xi + \alpha) - V_1(\xi) + \alpha_1(\xi) = \bar{\alpha}_1.
\]

Then the transformed map is
\[
\tilde{F}_1^\phi := S_1 \circ \bar{F}_1^\phi \circ S_1^{-1} : (\xi_1, \eta_1) \mapsto (\tilde{\xi}_1, \tilde{\eta}_1),
\]
with
\[
\left\{
\begin{array}{ll}
\tilde{\xi}_1 = \xi_1 + \alpha + \bar{\alpha}_1 \cdot (\lambda \eta_1) + \alpha_2(\xi_1)(\lambda \eta_1)^2 + O(\varepsilon^\gamma |\lambda \eta_1|^3) + O(\varepsilon^\gamma |\nu|), \\
\tilde{\eta}_1 = \nu + \bar{\beta}_1 \cdot (\lambda \eta_1) + \beta_2(\xi_1)(\lambda \eta_1)^2 + O(\varepsilon^\gamma |\lambda \eta_1|^3) + O(\varepsilon^\gamma |\nu|).
\end{array}
\right.
\]

\subsubsection{Step 2: Elimination of the non-constant quadratic term}

Let $\bar{\beta}_2 := \int_\mathbb{T} \beta_2(\xi)\, d\xi$. Consider the cohomological equation
\[
\bar{\beta}_1^2 U_2(\xi + \alpha) - \bar{\beta}_1 U_2(\xi) + \beta_2(\xi) = \bar{\beta}_2,
\]
which admits a unique solution $U_2 \in C^\omega(\mathbb{T})$. Define the coordinate transformation $T_2: (\xi_1, \eta_1) \mapsto (\xi_1, \eta_2)$ by
\[
\eta_2 = \eta_1 + U_2(\xi_1) \eta_1^2.
\]
This change transforms the non-constant coefficient $\beta_2(\xi_1)$ into the constant $\bar{\beta}_2$.

Similarly, the equation
\[
V_2(\xi + \alpha) - V_2(\xi) + \alpha_2(\xi) = \bar{\alpha}_2
\]
admits a unique solution $V_2 \in C^\omega(\mathbb{T})$. Define the coordinate transformation $S_2: (\xi_1, \eta_2) \mapsto (\xi_2, \eta_2)$ by
\[
\xi_2 = \xi_1 + V_2(\xi_1) \eta_2^2,
\]
which transforms the non-constant coefficient $\alpha_2(\xi_1)$ into the constant $\bar{\alpha}_2$.

Let $\tilde{F}_2^\phi := S_2 \circ T_2 \circ \tilde{F}_1^\phi \circ T_2^{-1} \circ S_2^{-1}$ denote the resulting map in the new coordinates.

Finally, by taking $\tilde{s}$ to be the minimum among $\bar{s}$ and the analyticity radii of $U_i$ and $V_i$ (for $i = 1, 2$), we complete the proof of Lemma~\ref{ketra} by relabeling the coordinates and the map:
\[
\tilde{F}_2^\phi \rightsquigarrow \tilde{F}^\phi, \quad (\xi_2, \eta_2) \rightsquigarrow (X, Y).
\]


 \subsection{Reduction of the perturbation}



Denote the integrable part of $\tilde{F}^\phi(X,Y)$ by $N(X,Y)$, that is,
\[
N(X,Y) = \left(X + \alpha + \bar{\alpha}_1 \cdot (\lambda Y) + \bar{\alpha}_2 \cdot (\lambda Y)^2,\; \nu + \bar{\beta}_1 \cdot (\lambda Y) + \bar{\beta}_2 \cdot (\lambda Y)^2\right).
\]
Consider the fixed point equation
\[
Y = \nu + \bar{\beta}_1 \cdot (\lambda Y) + \bar{\beta}_2 \cdot (\lambda Y)^2,
\]
whose solutions are given explicitly by
\[
Y_{\pm} = \frac{-(\bar{\beta}_1 \lambda - 1) \pm \sqrt{(\bar{\beta}_1 \lambda - 1)^2 - 4\nu \bar{\beta}_2 \lambda^2}}{2 \bar{\beta}_2 \lambda^2}.
\]
Recall that $\varepsilon = 1 - \lambda$. Then we have
\[
\bar{\beta}_1 \lambda - 1 = O(\varepsilon), \quad 4\nu \bar{\beta}_2 \lambda^2 = O(\varepsilon^{2\gamma}),
\]
which implies
\[
Y_{\pm} = O(\varepsilon^{\gamma - 1}).
\]

We now perform the final coordinate change $W : (X, Y) \mapsto (X, Z)$ defined by
\[
Z = Y - Y_+.
\]
Then the transformed map becomes
\[
\hat{F}^\phi := W \circ \tilde{F}^\phi \circ W^{-1} : (X, Z) \mapsto (\hat{X}, \hat{Z}).
\]
Using the identity
\[
Y_+ = \nu + \bar{\beta}_1 \cdot (\lambda Y_+) + \bar{\beta}_2 \cdot (\lambda Y_+)^2,
\]
we obtain the expansion
\[
\left\{
\begin{array}{ll}
\hat{X} = X + \alpha + \bar{\alpha}_1 \cdot (\lambda Y_+) + (1 + O(\varepsilon^\gamma)) \cdot (\lambda Z) + O(\varepsilon^\gamma Y_+^2) + O(\varepsilon^{2\gamma}), \\
\hat{Z} = (1 + O(\varepsilon^\gamma)) \cdot (\lambda Z) + O(\varepsilon^\gamma |Y_+|^3) + O(\varepsilon^{2\gamma}).
\end{array}
\right.
\]

\begin{Remark}
We could also perform the change of coordinates $Z = Y - Y_-$ instead of $Y_+$. Since $\gamma \in (1,2)$, the mixed and higher-order terms such as $O(\varepsilon^\gamma Y_+ Z)$, $O(\varepsilon^\gamma Z^2)$, and $O(\varepsilon^\gamma Y_+^2 Z)$ are absorbed into $O(\varepsilon^\gamma Z)$ and $O(\varepsilon^{2\gamma})$, respectively.
\end{Remark}

Note that the quantities $\bar{\alpha}_1 \cdot (\lambda Y_+)$, $O(\varepsilon^\gamma Y_+^2)$, and $O(\varepsilon^\gamma |Y_+|^3)$ are constants. Define
\[
\hat{\alpha}_1 := \alpha + \bar{\alpha}_1 \cdot (\lambda Y_+) + O(\varepsilon^\gamma Y_+^2), \quad
\hat{\alpha}_2 := O(\varepsilon^\gamma |Y_+|^3).
\]
Then the map $\hat{F}^\phi$ can be written in the form
\[
\hat{F}^\phi(X, Z) = \left(X + \hat{\alpha}_1 + \lambda_1 Z + \hat{\phi}_1(X),\; \hat{\alpha}_2 + \lambda_2 Z + \hat{\phi}_2(X)\right),
\]
where
\[
\lambda_1, \lambda_2 = (1 + O(\varepsilon^\gamma)) \cdot (1 - \varepsilon) = 1 - \varepsilon + O(\varepsilon^\gamma) + O(\varepsilon^{1 + \gamma}) < 1,
\]
and
\[
\|\hat{\phi}_i\|_{C^1} = O(\varepsilon^{2\gamma}), \quad i = 1, 2.
\]

However, in general we cannot guarantee that $\lambda_1 = \lambda_2$ or $\hat{\phi}_1 = \hat{\phi}_2$. Therefore, we invoke Remark~\ref{needb} here. Observe that
\[
O(\varepsilon^{2\gamma}) <O(\varepsilon^2) = \left(1 - \sqrt{\tilde{\lambda}}\right)^2,
\]
where $\tilde{\lambda} := \max\{\lambda_1, \lambda_2\}$.

By applying Theorems~\ref{M2} and \ref{M22} to the map $\hat{F}^\phi$, we conclude that it admits a unique $C^1$ invariant graph. This completes the proof of Theorem~\ref{M222}.







 \vspace{2em}

	\section{\sc The regularity of invariant graphs}\label{S5}
We provide a proof of Theorem~\ref{M3}. To this end, we construct an explicit example of a Lipschitz perturbation satisfying condition~(\ref{contrac}) such that the corresponding invariant Lipschitz graph is non-differentiable at certain points. This demonstrates that Theorems~\ref{M2} and~\ref{M22} do not imply each other.


\subsection{On the Denjoy counterexample}
Following \cite[Chapter X]{H11} (see also \cite{Ar11}), we construct a $C^1$ Denjoy counterexample $\bar{g}:\T\to\T$ as follows.

Let $\omega \in \mathbb{R}\setminus\mathbb{Q}$ be an irrational number, $\eps > 0$ a fixed parameter, and $N \gg 1$ a sufficiently large constant. Define a sequence $(\ell_k)_{k\in\mathbb{Z}}$ by:
\begin{equation}\label{l_k}
\ell_k := \frac{a_N}{(|k| + N)(\log(|k| + N))^{1+\eps}},
\end{equation}
where the normalization constant $a_N > 0$ is chosen so that $\sum_{k \in \mathbb{Z}} \ell_k = 1$.

Fix a $C^\infty$ bump function $\eta : \mathbb{R} \to \mathbb{R}$ satisfying:
\begin{itemize}
    \item $\eta \geq 0$,
    \item $\operatorname{supp}(\eta) \subset \left[\frac{1}{4}, \frac{3}{4}\right]$,
    \item $\int_0^1 \eta(t)\,dt = 1$.
\end{itemize}
For each $k \in \mathbb{Z}$, define the rescaled function:
\[
\eta_k(t) := \eta\left(\frac{t}{\ell_k}\right).
\]
Then $\int_0^{\ell_k} \eta_k(t)\,dt = \ell_k$, and there exist constants $C_1, C_2 > 0$ (depending only on $\eta$) such that:
\[
C_1 \leq \|\eta_k\|_{C^0} \leq C_2, \quad \text{and} \quad \frac{C_1}{\ell_k} \leq \|\eta_k'\|_{C^0} \leq \frac{C_2}{\ell_k}.
\]
From (\ref{l_k}), it follows that if $N$ is sufficiently large, then for all $k \in \mathbb{Z}$:
\begin{equation}\label{del_k}
\frac{C_1}{|k| + N} \leq \left|\frac{\ell_{k+1}}{\ell_k} - 1\right| \leq \frac{C_2}{|k| + N}.
\end{equation}

Now define a family of $C^\infty$ diffeomorphisms $b_k : [0, \ell_k] \to [0, \ell_{k+1}]$ by:
\[
\zeta_k(x) := \int_0^x \left[1 + \left(\frac{\ell_{k+1}}{\ell_k} - 1\right)\eta_k(t)\right] dt,
\]
so that $\zeta_k(\ell_k) = \ell_{k+1}$.

There exists a Cantor set $\mathcal{K} \subset \mathbb{T} := \mathbb{R}/\mathbb{Z}$ with $\operatorname{Leb}(\mathcal{K}) = 0$, whose complementary intervals $\{I_k\}_{k\in\mathbb{Z}}$ satisfy:
\begin{itemize}
    \item The ordering of $\{I_k\}$ on $\mathbb{T}$ follows the sequence $\{k\omega \mod 1\}_{k\in\mathbb{Z}}$;
    \item $\operatorname{length}(I_k) = \ell_k$ for each $k \in \mathbb{Z}$.
\end{itemize}

Define a semiconjugacy $j : \mathbb{T} \to \mathbb{T}$ to the rigid rotation $R_\omega(x) = x + \omega \mod 1$ via the probability measure:
\[
\bar{\mu} := \sum_{k\in\mathbb{Z}} \ell_k\, \delta_{k\omega},
\]
where $\delta_{k\omega}$ denotes the Dirac mass at $k\omega$. Then for $x \in \{k\omega\}_{k\in\mathbb{Z}}$, define:
\[
j^{-1}(x) := \int_0^x d\bar{\mu}(t).
\]

For each $I_k$, there exists $\lambda_k \in \mathbb{R}$ such that $R_{\lambda_k}(I_k) = [0, \ell_k]$. Define:
\[
\bar{g}_k := R_{\lambda_{k+1}} \circ \zeta_k \circ R_{-\lambda_k},
\]
so that $\bar{g}_k : I_k \to I_{k+1}$ is a $C^\infty$ diffeomorphism.

By \cite[Chapter X, (3.12)¨C(3.13)]{H11}, there exists a $C^1$ diffeomorphism $\bar{g} : \mathbb{T} \to \mathbb{T}$ such that:
\begin{itemize}
    \item $\bar{g}|_{I_k} = \bar{g}_k : I_k \to I_{k+1}$ (wandering intervals);
    \item $j \circ \bar{g} = R_\omega \circ j$ (semiconjugacy).
\end{itemize}
As a consequence:
\begin{itemize}
    \item $\bar{g}(\mathcal{K}) = \mathcal{K}$ (minimal invariant set);
    \item $\rho(\bar{g}) = \omega$ (rotation number).
\end{itemize}

By definition, the derivative of $\bar{g}$ on $I_k$ is:
\[
D\bar{g}_k = D\bar{g}|_{I_k} = \left( 1 + \left( \frac{\ell_{k+1}}{\ell_k} - 1 \right) \eta_k \right) \circ R_{-\lambda_k}.
\]
Furthermore, we obtain the following derivative estimates:
\begin{equation}\label{gkgkp}
\lim_{|k|\to\infty} \|D\bar{g}_k - 1\|_{C^0} = 0, \quad \text{and} \quad D\bar{g}(x) = 1, \quad \forall x \in \mathcal{K}.
\end{equation}

We denote by $g:\mathbb{R} \to \mathbb{R}$ the lift of $\bar{g} : \mathbb{T} \to \mathbb{T}$.

\subsection{Arnaud's modification}



Inspired by Arnaud \cite{Ar11}, we modify the Denjoy counterexample along a certain wandering orbit such that the resulting map $\bar{g}$ becomes non-differentiable at each point of that orbit. At the same time, we control the Lipschitz semi-norm of the perturbation $\phi : \T \to \R$ to ensure
\[
\|\phi\|_{\mathrm{Lip}} \leq (1 - \sqrt{\lambda})^2,
\]
where $\phi$ is defined as in Proposition \ref{hf1}. More specifically, let
\[
\varphi(x) := g(x) + \lambda g^{-1}(x) - (1 + \lambda)x,
\]
then set
\[
\phi(x) := \varphi(x) - \int_{\T} \varphi(x)\, dx.
\]
Consequently, the required dissipative twist map is given by
\begin{equation}\label{Fphi}
F^{\phi}(x, y) = \left(x + \alpha_1 + \lambda y + \phi(x),\ \lambda y + \phi(x)\right),
\end{equation}
where $\alpha_1 = \frac{1}{1 - \lambda} \int_{\T} \varphi(x)\, dx$.

Let $(I_k)_{k \in \mathbb{Z}}$ denote the family of connected components of $\mathbb{T} \setminus \mathcal{K}$, as defined previously. Fix a base point $x_0 \in I_0$ and consider its full orbit under $\bar{g}$:
\[
(x_k)_{k \in \mathbb{Z}} := (\bar{g}^k(x_0))_{k \in \mathbb{Z}}.
\]
Our goal is to construct a Lipschitz perturbation $\bar{h}$ of $\bar{g}$ such that:
\begin{enumerate}
    \item $\bar{h}|_{\mathcal{K}} \equiv \bar{g}|_{\mathcal{K}}$;
    \item $\bar{h}(x_k) = \bar{g}(x_k) = x_{k+1}$ for all $k \in \mathbb{Z}$.
\end{enumerate}
Let $h : \R \to \R$ denote the lift of $\bar{h} : \T \to \T$. For each $k \in \mathbb{Z}$, we introduce the following notation:
\begin{itemize}
    \item Interval decomposition:
    \[
    I_k := (a_k, b_k), \quad L_k := (a_k, x_k], \quad R_k := [x_k, b_k).
    \]
    \item Auxiliary function:
    \[
    W(x) := g(x) + \lambda g^{-1}(x) - (1 + \lambda)x.
    \]
    \item Derivative quantity:
    \[
    m_k := 1 + \lambda + DW(x_k),
    \]
    where $DW(x_k) := Dg(x_k) + \lambda Dg^{-1}(x_k) - (1 + \lambda)$. This implies the relation:
    \[
    Dg(x_k) + \frac{\lambda}{Dg(x_{k-1})} = m_k.
    \]
\end{itemize}

For each parameter $m > 2\sqrt{\lambda}$, define the M\"{o}bius-type transformation
\[
\Phi_m : (0, +\infty) \to (-\infty, m), \quad t \mapsto m - \frac{\lambda}{t},
\]
and denote by $\Phi_m^n(t)$ the $n$-fold composition of $\Phi_m$.

\begin{Proposition}[Properties of $\Phi_m$]\label{assyp}
For $m > 2\sqrt{\lambda}$, the map $\Phi_m$ satisfies the following:
\begin{enumerate}
    \item $\Phi_m$ is strictly increasing and bijective.
    \item When $m = 1 + \lambda$:
    \begin{itemize}
        \item The map has two fixed points: $p_- = \lambda$, $p_+ = 1$;
        \item If $t < \lambda$, then $\Phi_m^n(t) \to -\infty$ as $n \to +\infty$;
        \item If $t > \lambda$, then $\Phi_m^n(t) \to 1$ as $n \to +\infty$.
    \end{itemize}
    \item When $m \neq 1 + \lambda$:
    \begin{itemize}
        \item The map has two fixed points: $p_\pm = \frac{m \pm \sqrt{m^2 - 4\lambda}}{2}$;
        \item If $t < p_-$, then $\Phi_m^n(t) \to -\infty$ as $n \to +\infty$;
        \item If $t > p_-$, then $\Phi_m^n(t) \to p_+$ as $n \to +\infty$.
    \end{itemize}
\end{enumerate}
\end{Proposition}

By item (a), the inverse $\Phi^{-1}_m$ is well-defined and strictly increasing.

\begin{Proposition}[Properties of $\Phi^{-1}_m$]
For $m > 2\sqrt{\lambda}$, the inverse map $\Phi^{-1}_m$ satisfies:
\begin{enumerate}
    \item When $m = 1 + \lambda$:
    \begin{itemize}
        \item Fixed points: $p^m_- = \lambda$, $p^m_+ = 1$;
        \item If $t < 1$, then $\Phi_m^n(t) \to \lambda$ as $n \to -\infty$;
        \item If $t > 1$, then $\Phi_m^n(t) \to +\infty$ as $n \to -\infty$.
    \end{itemize}
    \item When $m \neq 1 + \lambda$:
    \begin{itemize}
        \item Fixed points: $p^m_\pm = \frac{m \pm \sqrt{m^2 - 4\lambda}}{2}$;
        \item If $t < p^m_+$, then $\Phi_m^n(t) \to \lambda$ as $n \to -\infty$;
        \item If $t > p^m_+$, then $\Phi_m^n(t) \to +\infty$ as $n \to -\infty$.
    \end{itemize}
\end{enumerate}
\end{Proposition}

Select an initial value $\beta_0 \neq \alpha_0$ (see Lemma \ref{forwap} below). Define the two-sided sequence $(\beta_k)_{k \in \mathbb{Z}}$ by the recurrence:
\[
\beta_k :=
\begin{cases}
\Phi_{m_k}(\beta_{k-1}), & k \geq 1, \\
\Phi^{-1}_{m_{k+1}}(\beta_{k+1}), & k \leq -1.
\end{cases}
\]

Recall the notation $N$ in the definition of $\ell_k$ (see (\ref{l_k})). Fix $\lambda \in (0,1)$ and define
\[
n_{\pm} := 1 + \lambda \pm \frac{1}{N}.
\]
For sufficiently large $N$, we ensure that for all $k \in \mathbb{Z}$:
\begin{equation}\label{mppp}
|m_k - (1 + \lambda)| \leq \frac{1}{N}, \quad
|p^{n_{\pm}}_+ - 1| \leq \frac{1}{\sqrt{N}}, \quad
|p^{n_{\pm}}_- - \lambda| \leq \frac{1}{\sqrt{N}}.
\end{equation}
Additionally, we impose the condition
\begin{equation}\label{lacond}
\frac{1}{\sqrt{N}} < \min\left\{\lambda,\ \frac{1}{2}(\sqrt{\lambda} - \lambda)\right\},
\end{equation}
which implies
\[
N > \max\left\{\frac{1}{\lambda^2},\ \frac{4}{\sqrt{\lambda}(1 - \sqrt{\lambda})^2}\right\}, \quad
\lambda + \frac{1}{\sqrt{N}} < \sqrt{\lambda} < 1 - \frac{1}{\sqrt{N}}.
\]
As $\lambda \to 0^+$, it suffices to choose $N \geq \lambda^{-2 - \varepsilon}$. As $\lambda \to 1^-$, it is enough to take $N \geq (1 - \lambda)^{-2 - \varepsilon}$.


\begin{Lemma}[Asymptotic behavior]\label{forwap}
Let $N$ satisfy conditions \eqref{del_k}, \eqref{mppp}, and \eqref{lacond}. Given $\lambda \in (0,1)$, for each
\begin{equation}\label{beN}
\beta_0 \in \left( \lambda + \frac{1}{\sqrt{N}},\; 1 - \frac{1}{\sqrt{N}} \right),
\end{equation}
the following limits hold:
\[
\lim_{k \to +\infty} \beta_k = 1, \quad \lim_{k \to -\infty} \beta_k = \lambda.
\]
\end{Lemma}

\begin{Remark}
By \eqref{mppp} and \eqref{beN}, we have $\beta_0 < \alpha_0$.
\end{Remark}

\begin{proof}
We prove only that $\lim_{k \to +\infty} \beta_k = 1$, since the case $\lim_{k \to -\infty} \beta_k = \lambda$ follows similarly by considering the inverse map $\Phi_m^{-1}$. The argument is inspired by \cite[Lemma 1]{Ar11}.

For all $k \geq 0$, we observe:
\[
\beta_{k+1} = \Phi_{m_{k+1}}(\beta_k) = \Phi_{n_+}(\beta_k) + m_{k+1} - n_+ < \Phi_{n_+}(\beta_k),
\]
and similarly,
\[
\beta_{k+1} = \Phi_{n_-}(\beta_k) + m_{k+1} - n_- > \Phi_{n_-}(\beta_k).
\]
Thus, for all $k \geq 0$, we obtain the inequality:
\begin{equation}\label{leqq}
\Phi_{n_-}^k(\beta_0) \leq \beta_k \leq \Phi_{n_+}^k(\beta_0).
\end{equation}

Since $\beta_0 > p_-^{n_-}$ by \eqref{mppp} and \eqref{beN}, Proposition~\ref{assyp} implies that
\[
\Phi_{n_-}^k(\beta_0) \to p_+^{n_-} \quad \text{as } k \to +\infty.
\]
Therefore, there exists $K_1$ such that for all $k \geq K_1$,
\[
|\Phi_{n_-}^k(\beta_0) - 1| \leq \frac{1}{\sqrt{N}},
\]
which in turn implies
\begin{equation}\label{bkpln}
\beta_k > p_-^{n_-}.
\end{equation}

For any fixed $\delta > 0$, define
\[
m_+ := 1 + \lambda + \delta, \quad m_- := 1 + \lambda - \delta.
\]
Taking $\delta > 0$ sufficiently small ensures that
\[
n_- < m_- < m_+ < n_+.
\]
This yields the comparison:
\begin{equation}\label{pcompar}
p_-^{n_+} < p_-^{m_+} < p_-^{m_-} < p_-^{n_-} < p_+^{n_-} < p_+^{m_-} < p_+^{m_+} < p_+^{n_+}.
\end{equation}

By the definition of $m_k$, there exists $K_2$ such that for all $k \geq K_2$,
\[
|m_k - (1 + \lambda)| \leq \delta.
\]
Arguing as in \eqref{leqq}, we get for all $K \geq 0$,
\[
\Phi_{m_-}^K(\beta_k) \leq \beta_{k+K} \leq \Phi_{m_+}^K(\beta_k).
\]
Let $K^* := \max\{K_1, K_2\}$. Then for $k \geq K^*$, it follows from \eqref{bkpln} and \eqref{pcompar} that
\[
\beta_{K^*} > p_-^{m_-} > p_-^{m_+}.
\]
By Proposition~\ref{assyp}, we conclude
\[
\Phi_{m_\pm}^K(\beta_{K^*}) \to p_+^{m_\pm} \quad \text{as } K \to +\infty.
\]
Moreover, by definition,
\[
p_+^{m_\pm} \to 1 \quad \text{as } \delta \to 0^+.
\]
Therefore, we obtain the desired limit:
\[
\lim_{k \to +\infty} \beta_k = 1.
\]
This completes the proof of Lemma \ref{beN}.\end{proof}

\medskip

Let $N$ satisfy conditions \eqref{del_k}, \eqref{mppp}, and \eqref{lacond}. We now construct $h_N$ by modifying $g_N$ so that $h_N$ fails to be differentiable at each point in the wandering orbit $\{x_k\}_{k \in \mathbb{Z}}$. To this end, it suffices to modify only the right-hand derivative of $g_N$ at $x_k$. Specifically, we set:
\[
D\bar{h}_N^R(x_k) = \beta_k, \quad D\bar{h}_N^L(x_k) = Dg_N(x_k).
\]
For $x_k \in (a_k, b_k)$, define the midpoint
\[
c_k := \frac{x_k + b_k}{2},
\]
and let
\[
J_k := (x_k, c_k), \quad J := \bigcup_{k \in \mathbb{Z}} J_k.
\]
It is clear that the Cantor set $\mathcal{K} \subset J$. Recall that $h_N : \mathbb{R} \to \mathbb{R}$ is the lift of $\bar{h}_N : \mathbb{T} \to \mathbb{T}$. We define $\bar{h}_N$ as follows:
\begin{itemize}
\item For $x \in \mathbb{T} \setminus J$, set $\bar{h}_N(x) = \bar{g}_N(x)$.
\item For $x \in J_k$, define $\bar{h}_N$ to be of class $C^1$ with derivative $D\bar{h}_N(x)$ lying between $\beta_k$ and $D\bar{g}_N(c_k)$.
\end{itemize}
It is straightforward to verify that $h_N \in D^0(\mathbb{T})$.






\subsection{Construction of the perturbation}


Consider the dissipative twist map defined by
\[
F^{\psi_N}(x,y) = \left(x + \lambda y + \psi_N(x),\; \lambda y + \psi_N(x)\right).
\]
By Proposition~\ref{hf1}, the map $F^{\psi_N}$ admits an invariant graph $\Psi_N$ if and only if
\[
\psi_N(x) = h_N(x) + \lambda h_N^{-1}(x) - (1 + \lambda)x,
\]
where $h_N \in D^0(\T)$ takes the form
\begin{equation}\label{hlpss}
h_N(x) = x + \lambda \Psi_N(x) + \psi_N(x).
\end{equation}
Let us define
\[
\psi_N(x) := h_N(x) + \lambda h_N^{-1}(x) - (1 + \lambda)x.
\]
Note that
\[
Dg_N\big|_{\T} = D\bar{g}_N\big|_{\T}, \quad Dh_N\big|_{\T} = D\bar{h}_N\big|_{\T}.
\]
By the definition of $\beta_k$, we have
\[
D\psi^L_N(x_k) + (1 + \lambda) = Dg_N(x_k) + \frac{\lambda}{Dg_N(x_{k-1})} = \beta_k + \frac{\lambda}{\beta_{k-1}} = D\psi^R_N(x_k) + (1 + \lambda).
\]
It follows that $\psi_N$ is differentiable at $x_k$ for all $k \in \mathbb{Z}$.

By definition, if $N$ is sufficiently large, then for all $k \in \mathbb{Z}$,
\[
\lambda - \frac{1}{\sqrt{N}} \leq p^{n_-}_- \leq \beta_k \leq p^{n_+}_- \leq 1 + \frac{1}{\sqrt{N}}, \quad \|\bar{g}_N - \mathrm{Id}\|_{C^0} \leq \frac{1}{N}.
\]
It follows from the construction of $\bar{h}_N$ that for all $x \in J$,
\[
\min_{k \in \mathbb{Z}} \left\{ \beta_k,\; D\bar{g}_N(c_k) \right\} \leq Dh_N(x) \leq \max_{k \in \mathbb{Z}} \left\{ \beta_k,\; D\bar{g}_N(c_k) \right\},
\]
which, together with the identity $\bar{h}_N(x) = \bar{g}_N(x)$ for $x \in \T \setminus J$, implies that for all $x \in \T$,
\[
\lambda - \frac{1}{\sqrt{N}} \leq Dh_N(x) \leq 1 + \frac{1}{\sqrt{N}}.
\]
This gives rise to the estimate for the Lipschitz semi-norm:
\begin{equation}\label{lipsem}
\|\psi_N\|_{\mathrm{Lip}} = \|D\psi_N\|_{L^\infty} \leq \left| \sqrt{\lambda} + \frac{\lambda}{\sqrt{\lambda}} - (1 + \lambda) \right| = (1 - \sqrt{\lambda})^2.
\end{equation}

Note that for each $k \in \mathbb{Z}$, the function $\psi_N$ is of class $C^1$ on $(a_k, b_k) \setminus \{x_k\}$. Meanwhile, the derivative of $\psi_N$ is uniformly bounded on $(a_k, b_k)$ for all $k \in \mathbb{Z}$. It follows that $\psi_N$ is of class $C^1$ on $\T \setminus \mathcal{K}$. In view of~\eqref{hlpss}, we observe that the graph $\Psi_N$ is not differentiable at $x_k$ for all $k \in \mathbb{Z}$. Hence, by Theorem~\ref{M22}, $\psi_N$ is not of class $C^1$. Moreover, $\psi_N$ must have non-differentiable points within the Cantor set $\mathcal{K}$.

Define
\[
\phi_N(x) := \psi_N(x) - \int_{\T} \psi_N(x)\, dx.
\]
The required dissipative twist map is then given by
\begin{equation}
F_{\alpha_N}^{\phi_N}(x,y) = \left(x + \alpha_N + \lambda y + \phi_N(x),\; \lambda y + \phi_N(x)\right),
\end{equation}
with
\[
\alpha_N = \frac{1}{1 - \lambda} \int_{\T} \psi_N(x)\, dx.
\]
Consequently, we prove Item (II) and Item (III) of Theorem~\ref{M3}. It remains to prove Item (I). To make the dependence of $\psi_N$ on $g$ or $h$ explicit, we define for $x \in \T$:
\[
\psi_N^g(x) := g_N(x) + \lambda g_N^{-1}(x) - (1 + \lambda)x, \quad A_N^g := \int_{\T} \psi_N^g(x) \, dx, \quad \phi_N^g(x) := \psi_N^g(x) - A_N^g,
\]
\[
\psi_N^h(x) := h_N(x) + \lambda h_N^{-1}(x) - (1 + \lambda)x, \quad A_N^h := \int_{\T} \psi_N^h(x) \, dx, \quad \phi_N^h(x) := \psi_N^h(x) - A_N^h.
\]

By the construction of $h_N$, the following properties hold:
\begin{itemize}
    \item If $x \in \mathcal{K}$, then $g_N(x) = h_N(x)$, which implies
    \begin{equation} \label{arg1}
        A_N^h + \phi_N^h(x) = \psi_N^h(x) = \psi_N^g(x) = A_N^g + \phi_N^g(x).
    \end{equation}

    \item For each $k \in \Z$, if $x \in (a_k, b_k)$, then
    \begin{equation} \label{arg2}
        A_N^h + \phi_N^h(x) = A_N^g + \phi_N^g(a_k) + \int_{a_k}^x D\phi_N^h(t) \, dt.
    \end{equation}
\end{itemize}

We now claim that $A_N^h = A_N^g$. Suppose, for contradiction, that
\[
\Delta := A_N^h - A_N^g > 0.
\]
Since
\[
\| Dg_N - 1 \|_{C^0} = O\left(\frac{1}{N}\right), \quad \text{as } N \to +\infty,
\]
it follows that $\| D\psi_N^g \|_{C^0} = O\left(\frac{1}{N}\right)$. Therefore,
\[
\| \phi_N^g \|_{C^0} \leq \| D\phi_N^g \|_{C^0} = \| D\psi_N^g \|_{C^0} = O\left(\frac{1}{N} \right).
\]

Moreover, we observe:
\begin{itemize}
    \item If $x \in \mathcal{K}$, then from \eqref{arg1}, $\phi_N^g(x) - \phi_N^h(x) = \Delta > 0$.
    \item If $x \in \T \setminus \mathcal{K}$, then there exists $k \in \Z$ such that $x \in (a_k, b_k)$, and by (\ref{arg2}), we have
    \[
    \phi_N^g(x) - \phi_N^h(x) = \Delta + \phi_N^g(x) - \phi_N^g(a_k) - \int_{a_k}^x D\phi_N^h(t) \, dt.
    \]
\end{itemize}

By the inequality~\eqref{lipsem}, we have $\| D\phi_N^h \|_{L^\infty} \leq (1 - \sqrt{\lambda})^2$. Combined with
\[
\| \phi_N^g \|_{C^0} = O\left( \frac{1}{N} \right), \quad x - a_k \leq b_k - a_k = O\left( \frac{1}{N} \right),
\]
we deduce that for $x \in \T \setminus \mathcal{K}$, it holds that $\phi_N^g(x) - \phi_N^h(x) > 0$. Hence,
\[
\int_{\T} \phi_N^g(x) \, dx > \int_{\T} \phi_N^h(x) \, dx,
\]
contradicting the fact that both integrals vanish. Therefore, we conclude that $A_N^h = A_N^g$.

With this equality, it follows from \eqref{arg1} and \eqref{arg2} that $\| \phi_N^h \|_{C^0} = O\left( \frac{1}{N} \right)$. Consequently, Item (I) follows directly from the interpolation inequality:
\[
\| \phi_N^h \|_{C^{1 - \varepsilon}}\lesssim \| \phi_N^h \|_{C^0}^{\varepsilon} \left( \| \phi_N^h \|_{C^0} + \| \phi_N^h \|_{\mathrm{Lip}} \right)^{1 - \varepsilon}.
\]

This completes the proof of Theorem~\ref{M3}.





\vspace{2em}




 \noindent\textbf{Data Availability Statement.}
The authors state that this manuscript has no associated data and there is no conflict of interest.


\medskip

\begin{thebibliography}{DGOW}
\renewcommand{\baselinestretch}{1.0}
\setlength\itemsep{-1pt}
%\def\bi#1{\bibitem[#1]{#1}}
\small

%\bibitem[Al]{A} J.Albrecht. {\it On the existence of invariant tori in nearly-integrable Hamiltonian systems with finitely
%differentiable perturbations}. Regular and Chaotic Dynamics.
%\textbf{12} (2007), 281-320.

\bibitem[AA24]{AA}  S. Allais, M.-C. Arnaud.
{\it The dynamics of conformal Hamiltonian flows: dissipativity and conservativity}. Rev. Mat. Iberoam. 40 (2024), no. 3, 987--1021.



\bibitem[Arn11]{Ar11}  M.-C. Arnaud.
{\it A non-differentiable essential irrational invariant curve
 for a $C^1$ symplectic twist map}. Journal of modern dynamics, 2011, 5 (3), pp.583-591.



\bibitem[AF24]{AF} M.-C. Arnaud, J. Fejoz.  {\it Invariant Submanifolds of conformal Symplectic Dynamics.}  Journal de l'\'{E}cole Polytechnique, 11 (2024), 159--185.


\bibitem[AHV24]{AHV} M.-C. Arnaud, V. Humili\`ere, C. Viterbo. {\it Higher Dimensional Birkhoff attractors}. ArXiv: 2404.00804, Prepint 2024.


\bibitem[AMS23]{ams} M.-C. Arnaud, J.E.  Massetti and A. Sorrentino.  {\it On the fragility of periodic tori for families of symplectic twist maps}.  Advances in Mathematics. {429} (2023), Published online.

\bibitem[Arn61]{Ar}  V. I. Arnold. {\it  Small denominators. I. Mapping the circle onto itself}.
Izv. Akad. Nauk SSSR Ser. Mat. 25 (1961), 21-86.

\bibitem[BB13]{BB} P. Berger  and A. Bounemoura. {\it A geometrical proof of the persistence of normally hyperbolic submanifolds}. Dynamical Systems, 2013, 28 (4), 567-581.

\bibitem[Bir32]{B} G. D. Birkhoff. {\it Sur quelques courbes ferm\'{e}es remarquables}. Bull. Soc. Math. France {60}
(1932), 1-26.



\bibitem[Boh84]{Bo} T. Bohr. {\it A bound for the existence of invariant circles in a class of two-dimensional dissipative maps}. Physics Letters A {104}
(1984), 441-443.


\bibitem[CCD13]{CCD1} R. Calleja, A. Celletti and R. de la Llave. {\it A KAM theory for conformally symplectic systems: efficient
algorithms and their validation}. J. Differ. Equ. {255} (2013), 978-1049.


\bibitem[CCD22]{CCD2} R. Calleja, A. Celletti and R. de la Llave. {\it KAM quasi-periodic solutions for the dissipative standard map}. Commun. Nonlinear Sci. Numer. Simul. 106 (2022), Paper No. 106111, 29 pp.


%\bibitem[CdS]{CdS}
%A. Cannas da Silva. {\it Lecture on Symplectic Geometry}. Lecture Notes in Mathematics, Vol. 1764, Springer-Verlag, Berlin (2001), xii+217 pp.


\bibitem[Cas87]{Ca} M. Casdagli. {\it Periodic orbits for dissipative twist maps}. Ergod. Th. \& Dynam. Sys. {7} (1987), 165-173.

%\bibitem[CW]{CW}
%C.-Q. Cheng and L. Wang. {\it Destruction of Lagrangian torus for
%positive definite Hamiltonian systems}. Geom. Funct. Anal. {23} (2013), 848-866.


\bibitem[Con78]{Co} C. Conley. {\it Isolated Invariant Sets and the Morse Index}. CBMS Regional Conference Series in Mathematics. {38} American Mathematical Society, Providence (1978).



\bibitem[dMvS93]{de}  W. de Melo, S. van Strien. {\it One-Dimensional Dynamics}. Springer-Verlag. (1993)


\bibitem[Den32]{De} A. Denjoy.
{\it Sur les courbes d\'{e}finies par les \'{e}quations diff\'{e}rentielles \`{a} la surface du tore}. J. de Math. Pures et Appl.,
 (9), 11 (1932), 333-375

%\bibitem[Dev]{Dev}  R. Devaney. {\it An Introduction to Chaotic Dynamical Systems}. Wetsview
% Press. 2 (2003), 102-113

%\bibitem[Fa]{Fav} J. Favard. {\it Sur les meilleurs proc\'{e}d\'{e}s d'approximation}. B.S.M. 61 (1937), 243-256.


\bibitem[For94]{F} G. Forni. {\it
Analytic destruction of invariant circles}. Ergod.
 Th. \& Dynam. Sys. {14} (1994), 267-298.


%\bibitem[Hal83]{Ha} G. Hall. {\it
%Bifurcation of an attracting invariant circle: a Denjoy attractor}. Ergod.
% Th. \& Dynam. Sys. {3} (1983), 87-118.


 \bibitem[Her79]{H11} M. R. Herman. {\it Sur la conjugation diff$\acute{e}$rentiable des diff$\acute{e}$omorphismes du cercle $\grave{a}$ des
 rotations}. Publ. Math. IHES {49} (1979), 5-233.


\bibitem[Her83]{H1} M.R. Herman. {\it Sur les courbes invariantes par les diff\'{e}omorphismes de
 l'anneau}. Ast\'{e}risque {103-104} (1983), 1-221.

\bibitem[Her86]{H33} M. R. Herman.  {\it Sur les courbes invariantes par les
diff$\acute{e}$omorphismes de
 l'anneau}. Ast$\acute{\text{e}}$risque {144} (1986), 1-243.

%\bibitem[H3]{H2} M.R. Herman. {\it In\'{e}galit\'{e}s ``a priori" pour des
%tores lagrangiens invariants par des diff\'{e}omrphismes
%symplectiques}. Inst. Hautes \'{E}tudes Sci. Publ. Math. {70}
%(1990), 47-101.

\bibitem[Her90]{H3} M.R. Herman. {\it Non existence of Lagrangian graphs}.
unpublished preprint (1990).


\bibitem[HPS77]{HPS} M. Hirsch, C. Pugh and M. Shub. {\it Invariant manifolds}. Lecture Notes in Mathematics, Vol. 583. Springer-Verlag, (1977), ii+149 pp.


\bibitem[LeC86]{L1} P.  Le Calvez. {\it Existence d'orbites quasi-periodiques dans les attracteurs de Birkhoff}.
Commun. Math. Phys. {106} (1986), 383-394.

\bibitem[LeC87]{L2} P.  Le Calvez. {\it Propri\'{e}t\'{e}s des attracteurs de Birkhoff}.
Ergod. Th. \& Dynam. Sys. {8} (1987), 241-310.


%\bibitem[Li]{Li} P. Libermann. {\it Sur les automorphismes infinit\'{e}simaux des structures symplectiques et des structures de contact}. (French) 1959 Colloque G\'{e}om. Diff. Globale
%(Bruxelles, 1958) pp. 37-59 Centre Belge Rech. Math., Louvain.


%
 \bibitem[Mat84]{M1} J. N. Mather. {\it Non-existence of invariant circles}. Ergod.
 Th. \& Dynam. Sys. {4} (1984), 301-309.



 \bibitem[Mat88]{M4} J. N. Mather. {\it Destruction of invariant circles}. Ergod.
 Th. \& Dynam. Sys. {8} (1988), 199-214.



\bibitem[Mas18]{Mas1} {J. Massetti. {\it A normal form \`{a} la Moser for diffeomorphisms and a generalization of R\"{u}ssmann's translated curve theorem
 to higher dimensions}. Anal. PDE, 11(1):149-170, 2018..}


\bibitem[Mas23]{Ma} {J. Massetti. {\it Attractive invariant circles \`{a} la Chenciner}. Regul. Chaot. Dyn. 28, 447-467 (2023).}

\bibitem[MS17]{MS} S. Mar\`o and A. Sorrentino.
{\it Aubry-Mather theory for conformally symplectic systems.} Comm. Math. Phys.  {354} (2017), 775-808.


\bibitem[New04]{Ne} S. Newhouse. {\it Cone-fields, domination, and hyperbolicity}. Brin,
 Michael (ed.) et al., Modern dynamical systems and applications.
 Dedicated to Anatole Katok on his 60th birthday. Cambridge University Press. 419-432, 2004.

%\bibitem[P]{P} J. P\"{o}schel.  {\it \"{U}ber invariante Tori in differenzierbaren Hamiltonschen Systemen}. (German)[On invariant tori in differentiable Hamiltonian systems]
%Bonner Math. Schriften, 120[Bonn Mathematical Publications]
%Universit\"{a}t Bonn, Mathematisches Institut, Bonn, 1980. 103 pp.



			\bibitem[R\"{u}s70]{R1} H. R\"{u}ssmann.  {\it Kleine Nenner. I. \"{U}ber invariante Kurven differenzierbarer Abbildungen eines Kreisringes}. Nachr. Akad. Wiss. G\"{o}ttingen Math.-Phys. Kl. II(1970), 67-105.



\bibitem[Sal04]{Sa} D. Salamon. {\it The Kolmogorov-Arnold-Moser theorem}.
Math. Phys. Electron. J. 10 (2004), Paper 3, 37 pp.

\bibitem[SW25]{SoW} A. Sorrentino and L. Wang. {\it On the Destruction of Invariant Lagrangian Graphs for Conformal Symplectic Twist Maps}, {arXiv}, 2025.

%\bibitem[SuW]{SW} X. Su and L. Wang. {\it Total destruction of invariant tori for the generalized Frenkel-Kontorova model}, {Journal of Math. Phys.}, 53 (2012), 022702, 7pp.
%
%
%\bibitem[Ta]{T} F. Takens. {\it A $C^1$ counterexample to Moser's twist theorem. } Nederl. Akad. Wetensch. Proc. Ser. A {74}. Indag. Math.  {33}
%(1971), 378-386.
%
%
%	\bibitem[TL]{TL} {Z. Tong and Y. Li. {\it Towards continuity: Universal frequency-preserving KAM persistence and remaining regularity in nearly non-integrable Hamiltonian systems}.
%			Communications in Contemporary Mathematics. Published online (2024).}


\bibitem[Wan12]{W} L. Wang. {\it Variational destruction of invariant
circles}. {Discrete Contin. Dyn. Syst.}, 32
(2012), 4429-4443.

%\bibitem[W2]{W2}L. Wang. {\it Total destruction of Lagrangian tori}. J. Math. Anal. Appl. {410} (2014), 827-836.012).

\bibitem[Yoc84]{Y1} J.-C. Yoccoz. {\it
Conjugaison diff\'{e}rentiable des diff\'{e}omorphismes du cercle dont le nombre de rotation v\'{e}rifie une condition diophantienne}.
Ann. Sci. \'{E}cole Norm. Sup. (4) 17 (1984), no. 3, 333-359.













%\bibitem[Y2]{Y2} J.-C. Yoccoz. {\it
%Centralisateurs et conjugaison diff\'{e}rentiable des diff\'{e}omorphismes du cercle}.
%Ast\'{e}risque No. 231 (1995), 89-242.
%
%\bibitem[Y3]{Y3} J.-C. Yoccoz. {\it
%Analytic linearization of circle diffeomorphisms}.
%Lecture Notes in Math., 1784
%Fond. CIME/CIME Found. Subser.
%Springer-Verlag, Berlin, 2002, 125-173.




%\bibitem[Za]{Za} G. Zaslavsky. {\it The simplest case of a strange attractor. } Physics Letters. {69A}
%(1978), 145-147.



%\bibitem[Zy]{Z} A. Zygmund. {\it Trigonometric Series}. Third Edition Volumes I \& II combined, with a foreword by Robert Fefferman. Cambridge
%University Press, Cambridge, 2002.
\end{thebibliography}








\end{document}


}
